\providecommand{\anapos}[1]{\bigl(#1\bigr)_{+}}
\providecommand{\ananeg}[1]{\bigl(#1\bigr)_{-}}
\providecommand{\anaone}{\mathbf{1}}
\providecommand{\anaE}{\mathcal{E}}
\providecommand{\anaEx}{\mathbb{E}}
\providecommand{\anaPr}{\mathbb{P}}

\section{One-sided analytic criteria}\label{sec:ana}

The source theorem is obtained from an upper bound for a smoothed sum over primes. A Mellin
transform turns that bound into a statement about poles, and a positivity argument of Landau
excludes them (\cref{src:sec:mechanism}). This section states that analytic step once, in the
generality needed in \cref{sec:nab,sec:grp,sec:hec}, together with its variants: tests in mean,
tests without a centring constant, joint tests, the removal of prime powers, the choice of the
auxiliary second scale, and tests on thin sets of twists and of scales. With the exceptions
marked ``\stG'', the results are unconditional statements of analysis. None of them supplies
the arithmetic inequality that a criterion asks for. For every function outside the scope of
\cref{src:thm:dgrh} that inequality is open.

For a real number $u$ write $u_+=\max(u,0)$ and $u_-=\max(-u,0)$. Mellin transforms are taken
over $[1,\infty)$. With $x=e^v$ the integral $\int_1^\infty R(x)x^{-s}\,\dd x$ is the shifted
Laplace transform $\int_0^\infty R(e^v)e^{(1-s)v}\,\dd v$ used in the Lean source.

\subsection{The basic test}\label{ana:sec:basic}
{\raggedright\sources{grh_transfer_analysis; grh_transfer_analysis_2026-11-25; grh_lean_transfer_analysis; grh_method_transfers_2026-11-25_2; grh_transfer_synthesis_2026-11-25; grh_sylow_transfer_2026-11-25; grh_reflection_transfer_2026-11-25; grh_cyclic_generator_transfer_2026-11-25; grh_a6_transfer_2026-11-25; grh_kummer_window_transfer_2026-11-25; grh_relative_prime_races_2026-11-25; psl2_grh_transfer}}

\begin{lemma}[Landau]\label{ana:lem:landau}
Let $h\colon[1,\infty)\to[0,\infty)$ be locally integrable, and suppose that
$G(s)=\int_1^\infty h(x)x^{-s}\,\dd x$ converges for some real $s$. Let $a\in[-\infty,\infty)$
be its abscissa of convergence. If $a>-\infty$, then $G$ has no holomorphic continuation to a
neighbourhood of the real point $a$.
\end{lemma}

\begin{proof}
Suppose that $G$ is holomorphic on a disc about $a$. Choose a real $z>a$ so close to $a$ that
the Taylor disc of $G$ about $z$ contains a real point $t=z-r<a$. For $k\ge0$,
differentiation inside the half-plane of convergence gives
\[
 (-1)^kG^{(k)}(z)=\int_1^\infty(\log x)^kh(x)x^{-z}\,\dd x .
\]
All terms of the Taylor series of $G$ about $z$, evaluated at $t$, are therefore nonnegative,
and Tonelli's theorem identifies its convergent sum with
\[
 \sum_{k\ge0}\frac{r^k}{k!}\int_1^\infty(\log x)^kh(x)x^{-z}\,\dd x
 =\int_1^\infty h(x)x^{-t}\,\dd x .
\]
So the integral converges at $t<a$, which contradicts the definition of $a$. See
\cite[\S15.1]{MontgomeryVaughan2007Multiplicative}.
\end{proof}

\begin{theorem}[One-sided Mellin--Landau test; \stU]\label{ana:thm:onesided}
Let $\beta>0$, and let $R\colon[1,\infty)\to\R$ be locally integrable with $\abs{R(x)}\le Cx^{A}$
for all sufficiently large $x$, for some finite constants $A$ and $C$. Suppose that
\[
 F(s)=\int_1^\infty R(x)x^{-s}\,\dd x\qquad(\Re s>A+1)
\]
has a meromorphic continuation to $\Re s>\beta$, and that every pole of $F$ at a real point
$u>\beta$ is simple with positive residue. Poles of $F$ off the real axis may have any order.
If
\begin{equation}\label{ana:eq:posint}
 \int_1^\infty R(x)_+\,x^{-\sigma}\,\dd x<\infty\qquad\text{for every real }\sigma>\beta ,
\end{equation}
then $F$ is holomorphic on $\Re s>\beta$. Moreover
$\int_1^\infty\abs{R(x)}x^{-\sigma}\,\dd x<\infty$ for every $\sigma>\beta$, so the Mellin integral
of $R$ converges absolutely on that half-plane.
\end{theorem}

\begin{proof}
Put $H(s)=\int_1^\infty R(x)_+x^{-s}\,\dd x$. By \eqref{ana:eq:posint} this integral converges
absolutely and defines a holomorphic function on $\Re s>\beta$. Far to the right,
\[
 G(s):=H(s)-F(s)=\int_1^\infty R(x)_-\,x^{-s}\,\dd x
\]
is the Mellin transform of a nonnegative function. Let $a$ be its abscissa of convergence. The
growth hypothesis gives $a\le A+1$, and it is enough to prove $a\le\beta$.

Suppose $a>\beta$. The meromorphic function $H-F$ continues $G$ to a neighbourhood of $a$, and
agrees with $G$ on $\Re s>a$ by the identity theorem. By hypothesis its singularity at $a$ is at
most a simple pole, with residue $-\Res_{s=a}F\le0$. On the other hand $G(\sigma)\ge0$ for real
$\sigma>a$, so that residue equals $\lim_{\sigma\downarrow a}(\sigma-a)G(\sigma)\ge0$. Hence the
residue is zero and $H-F$ is holomorphic at $a$. This contradicts \cref{ana:lem:landau}, and so
$a\le\beta$.

The positive and negative parts of $R$ now both have absolutely convergent transforms at every
real $\sigma>\beta$. Their difference is holomorphic on $\Re s>\beta$ and equals $F$ by analytic
continuation. Their sum gives the last assertion.
\end{proof}

The Lean source proves the special case in which a fixed exponential majorant is subtracted
(\code{IsPolar.landau}). The proof above subtracts the transform of the whole positive part
instead, which is what admits hypotheses in mean. The positivity principle is classical, and
related one-sided criteria for the zeta function and for Dirichlet $L$-functions appear in
\cite{Suzuki2025Variants}.

\begin{corollary}[Pointwise, dyadic and square-mean hypotheses]\label{ana:cor:forms}
Under the analytic hypotheses of \cref{ana:thm:onesided}, each of the following conditions
implies \eqref{ana:eq:posint}, and hence excludes all poles of $F$ in $\Re s>\beta$.
\begin{enumerate}[label=\textup{(\alph*)}]
\item For every $\eps>0$ there is $C_\eps$ with $R(x)\le C_\eps x^{\beta-1+\eps}$ for all
 sufficiently large $x$. In particular $R(x)\le Cx^{\beta-1}$ eventually suffices.
\item For every $\eps>0$ and all sufficiently large dyadic $X=2^k$,
 $\int_X^{2X}R(x)_+\,\dd x\ll_\eps X^{\beta+\eps}$.
\item For every $\eps>0$ and all sufficiently large dyadic $X=2^k$,
 $\int_X^{2X}R(x)_+^2\,\dd x\ll_\eps X^{2\beta-1+\eps}$.
\end{enumerate}
Condition \textup{(b)} is equivalent to \eqref{ana:eq:posint}.
\end{corollary}

\begin{proof}
For $\sigma>\beta$ choose $0<\eps<\sigma-\beta$. Under (a) the integrand in
\eqref{ana:eq:posint} is eventually at most $C_\eps x^{\beta-1+\eps-\sigma}$. Under (b),
\[
 \int_{2^k}^{2^{k+1}}R(x)_+\,x^{-\sigma}\,\dd x\ll_\eps2^{k(\beta+\eps-\sigma)},
\]
and the geometric series converges. Conversely, integrability at $\sigma=\beta+\eps$ bounds a
dyadic block by a constant multiple of $X^{\beta+\eps}$. Finally the Cauchy--Schwarz inequality
gives $\int_X^{2X}R_+\le X^{1/2}\bigl(\int_X^{2X}R_+^2\bigr)^{1/2}$, which takes (c) to (b).
\end{proof}

When a real pole may have higher order, the pointwise form of the test still holds if one
inspects the leading Laurent coefficient.

\begin{proposition}[Real poles of higher order; \stU]\label{ana:prop:higher}
Let $a_0<\beta$ be real, and let $R\colon[1,\infty)\to\R$ be locally integrable and of at most
polynomial growth, with $R(x)\le Cx^{\beta-1}$ for all sufficiently large $x$. Suppose that the
Mellin transform $F$ of $R$ continues meromorphically to $\Re s>a_0$, and that at every real
pole $a>\beta$ its leading Laurent coefficient is positive: if the pole has order $j$, then
$F(s)=c_a(s-a)^{-j}+O\bigl((s-a)^{1-j}\bigr)$ with $c_a>0$. Then $F$ has no pole in
$\Re s>\beta$, and the Mellin integral of $R$ converges absolutely there.
\end{proposition}

\begin{proof}
Enlarge $C$ to be nonnegative and replace $R$ by zero on an initial compact interval, so that
the upper bound holds for all $x\ge1$. The change in $F$ is entire and affects no pole and no
leading coefficient. Put
\[
 h(x)=Cx^{\beta-1}-R(x)\ge0,\qquad
 G(s)=\int_1^\infty h(x)x^{-s}\,\dd x=\frac{C}{s-\beta}-F(s)
\]
in the initial domain of convergence, and let $\sigma$ be the abscissa of convergence of this
nonnegative integral. Suppose $\sigma>\beta$. By \cref{ana:lem:landau} the meromorphic
continuation of $G$ has a pole at $\sigma$. Since $C/(s-\beta)$ is holomorphic there, the
leading Laurent coefficient of $G$ at $\sigma$ is the negative of that of $F$, hence negative,
and $G(s)\to-\infty$ as real $s\downarrow\sigma$. But $G(s)\ge0$ for real $s>\sigma$. Thus
$\sigma\le\beta$, $G$ is holomorphic on $\Re s>\beta$, and $F$ has no pole there. Finally
$\abs{R(x)}\le Cx^{\beta-1}+h(x)$ gives the absolute convergence.
\end{proof}

\begin{remark}[The hypotheses cannot be omitted]\label{ana:rem:hyp}
Fix $0<\beta<\alpha<1$ and $\tau>0$. Each of the following functions is nonpositive on
$[1,\infty)$, so every form of the one-sided hypothesis holds, and each has a transform with a
pole to the right of $\beta$.
\begin{enumerate}[label=\textup{(\arabic*)}]
\item \emph{A negative real residue.} $R(x)=-x^{\alpha-1}$ has $F(s)=-1/(s-\alpha)$.
\item \emph{A real pole masking complex ones.} $R(x)=x^{\alpha-1}\bigl(\cos(\tau\log x)-1\bigr)$
 has
 \[
  F(s)=\frac{s-\alpha}{(s-\alpha)^2+\tau^2}-\frac1{s-\alpha},
 \]
 with poles at $\alpha\pm i\tau$ of residue $1/2$ and a real pole at $\alpha$ of residue $-1$.
\item \emph{Higher order at a real point.} $R(x)=-x^{\alpha-1}\log x$ has
 $F(s)=-(s-\alpha)^{-2}$. The residue is zero although a double pole exists. This is why the
 source proves that its poles are at most simple, and why \cref{ana:prop:higher} inspects the
 leading coefficient.
\end{enumerate}
There is a separate cancellation problem. A difference of logarithmic derivatives does not see
a zero shared by its two factors: for $F_1(s)=F_2(s)=s-\rho$ one has $F_1'/F_1-F_2'/F_2=0$. A
test applied to a remainder must therefore identify the residues and prevent cancellation. In
the settings of \cref{ana:sec:ramp} this is automatic, because all multiplicities enter with
one sign. In the source the residues are differences of multiplicities, and an auxiliary twist
without zeros in a prescribed box is selected first (\cref{ana:sec:sparse}).
\end{remark}

\subsection{Ramp sums of Dirichlet series}\label{ana:sec:ramp}
{\raggedright\sources{grh_sylow_transfer_2026-11-25; grh_reflection_transfer_2026-11-25; psl2_grh_transfer; grh_a6_transfer_2026-11-25; grh_relative_prime_races_2026-11-25; grh_cyclic_generator_transfer_2026-11-25; grh_deficit_transfer_2026-11-25; grh_kummer_window_transfer_2026-11-25; grh_transfer_analysis_2026-11-25}}

\begin{definition}[Ramp and centred remainder]\label{ana:def:ramp}
Let $a(n)$, $n\ge1$, be real numbers with $a(n)=O(\log 2n)$, and put
\[
 A(s)=\sum_{n\ge1}a(n)n^{-s}\quad(\Re s>1),\qquad
 T(x)=\sum_{n\le x}a(n)\Bigl(\frac1n-\frac1x\Bigr)\quad(x\ge1).
\]
Suppose that $A$ extends meromorphically to a neighbourhood of $s=1$ with at most a simple pole
there, and write $A(s)=\delta/(s-1)+b_0+O(s-1)$. The \emph{centred remainder} of $a$ is
\[
 R(x)=T(x)-\delta\log x-c,\qquad c=b_0-\delta .
\]
\end{definition}

The kernel $1/n-1/x$ is the one of the source's \code{primeError}. Each summand vanishes when it
enters, so $T$ is continuous.

\begin{lemma}[Mellin transform of the centred ramp]\label{ana:lem:mellin}
In the situation of \cref{ana:def:ramp}:
\begin{enumerate}[label=\textup{(\roman*)}]
\item $T(x)=\int_1^x\bigl(\sum_{n\le t}a(n)\bigr)t^{-2}\,\dd t$ and $T(x)=O(\log^2(2x))$.
\item For $\Re s>1$,
 \begin{equation}\label{ana:eq:mellin}
  \int_1^\infty T(x)x^{-s}\,\dd x=\frac{A(s)}{s(s-1)},\qquad
  \Phi(s):=\int_1^\infty R(x)x^{-s}\,\dd x
  =\frac{A(s)}{s(s-1)}-\frac{\delta}{(s-1)^2}-\frac{c}{s-1}.
 \end{equation}
\item $\Phi$ is regular at $s=1$, and $c$ is the only constant with this property. If
 $T(x)-\delta\log x$ has a finite limit as $x\to\infty$, that limit is $c$.
\item On any domain in $\Re s>0$ to which $A$ continues meromorphically, so does $\Phi$. A pole
 of $A$ at $\rho\ne1$ is a pole of $\Phi$ of the same order. If it is simple with residue $r$,
 then $\Res_{s=\rho}\Phi=r/\bigl(\rho(\rho-1)\bigr)$.
\end{enumerate}
\end{lemma}

\begin{proof}
(i) Use $\int_n^xt^{-2}\,\dd t=1/n-1/x$ and $\sum_{n\le t}\abs{a(n)}\ll t\log 2t$. (ii) Absolute
convergence permits termwise integration, and for every $n\ge1$
\[
 \int_n^\infty\Bigl(\frac1n-\frac1x\Bigr)x^{-s}\,\dd x=\frac{n^{-s}}{s(s-1)},\qquad
 \int_1^\infty x^{-s}\log x\,\dd x=\frac1{(s-1)^2}.
\]
(iii) Division of $A(s)=\delta/(s-1)+b_0+O(s-1)$ by $s(s-1)$ gives
$\delta/(s-1)^2+(b_0-\delta)/(s-1)+O(1)$, so the two subtractions in \eqref{ana:eq:mellin}
remove the singularity exactly when $c=b_0-\delta$. Suppose $T(x)-\delta\log x\to c'$. For real
$\sigma>1$ and every $X\ge1$,
\[
 (\sigma-1)\Bigl|\int_1^\infty\bigl(T(x)-\delta\log x-c'\bigr)x^{-\sigma}\,\dd x\Bigr|
 \le(\sigma-1)\int_1^X\abs{T(x)-\delta\log x-c'}\,\dd x+\sup_{x>X}\abs{T(x)-\delta\log x-c'} ,
\]
which tends to zero as $\sigma\downarrow1$ and then $X\to\infty$. The left side equals
$\abs{(\sigma-1)\Phi(\sigma)+c-c'}$, and $\Phi$ is regular at $1$, so $c'=c$. (iv) is immediate
from \eqref{ana:eq:mellin}.
\end{proof}

\begin{definition}[One-sided setting]\label{ana:def:setting}
Fix $0<\beta<1$. A sequence $a$ as in \cref{ana:def:ramp} is in a \emph{one-sided setting at
$\beta$} if, for $\Re s>1$,
\begin{equation}\label{ana:eq:setting}
 A(s)=-\sum_{j=1}^{r}w_j\frac{F_j'(s)}{F_j(s)}+H(s)
\end{equation}
with positive real weights $w_j$, where $H$ is holomorphic on $\Re s>\beta$ and every $F_j$ is
meromorphic on $\Re s>\beta$, holomorphic and nonvanishing at every point of $\Re s\ge1$ other
than $s=1$, and holomorphic at every point of $\beta<\Re s<1$.
\end{definition}

In such a setting $A$ is meromorphic on $\Re s>\beta$. Its only possible pole on $\Re s\ge1$ is
at $s=1$, where it is simple with residue $\delta=-\sum_jw_j\ord_{s=1}F_j$. Its other poles lie
at the zeros $\rho$ of the $F_j$ in $\beta<\Re s<1$, and are simple with residue
$-\sum_jw_j\ord_\rho F_j$.

\begin{theorem}[One-sided zero test for ramp sums; \stU]\label{ana:thm:zerotest}
Let $a$ be in a one-sided setting at $\beta$, with centred remainder $R$. If
\[
 \int_1^\infty R(x)_+\,x^{-\sigma}\,\dd x<\infty\qquad\text{for every real }\sigma>\beta,
\]
and in particular if $R(x)\le Cx^{\beta-1}$ for all sufficiently large $x$, then no $F_j$ has a
zero in $\Re s>\beta$. The Mellin integral of $R$ then converges absolutely on $\Re s>\beta$.
\end{theorem}

\begin{proof}
By \cref{ana:lem:mellin} the transform $\Phi$ of $R$ is meromorphic on $\Re s>\beta$ and regular
at every point of $\Re s\ge1$. Its poles in $\beta<\Re s<1$ are simple, with
\begin{equation}\label{ana:eq:residue}
 \Res_{s=\rho}\Phi(s)=\frac{\sum_jw_j\ord_\rho F_j}{\rho(1-\rho)} .
\end{equation}
At a real point $\rho\in(\beta,1)$ the denominator is positive and the numerator is a
nonnegative combination of nonnegative integers. So the real poles of $\Phi$ have positive
residue, and $R(x)=O(\log^2(2x))$. \Cref{ana:thm:onesided} shows that $\Phi$ is holomorphic on
$\Re s>\beta$. A zero of some $F_j$ at a point $\rho$ of this half-plane would make
\eqref{ana:eq:residue} nonzero, because all orders are nonnegative, the weights are positive
and $\rho(1-\rho)\ne0$. No cancellation is possible, even at nonreal zeros.
\end{proof}

The positivity of the real residues is the point of the definition. It comes from the
holomorphy of the $F_j$ in the strip and needs no information about the individual Artin
factors of the $F_j$. No zero-density estimate and no auxiliary twist is used.

\begin{remark}[Orientation]\label{ana:rem:sign}
The inequality cannot be reversed. If $A=+\sum_jw_jF_j'/F_j+H$, as happens for a weight that
counts a deficit of the primes detected by the $F_j$, the real residues of $\Phi$ are negative.
The valid one-sided hypothesis is then an upper bound for $-R$, which is a lower bound for the
deficit statistic. An upper bound for it detects nothing. In the same way a real exceptional
zero would escape a lower bound for the statistics of \cref{ana:ex:instances}. A geometric
construction must therefore produce its inequality with the orientation of
\cref{ana:thm:zerotest}.
\end{remark}

\begin{example}[Instances]\label{ana:ex:instances}
The following settings are used in \cref{sec:nab,sec:grp,sec:hec}. In each of them the inputs
are classical: Hecke's continuation and the nonvanishing on $\Re s\ge1$ of the $L$-functions of
finite-order characters, and the pole of a Dedekind zeta function at $1$
\cite{Neukirch1999Algebraic}.
\begin{enumerate}[label=\textup{(\alph*)}]
\item \emph{A Dedekind zeta function.} For a number field $K$ take $A=-\zeta_K'/\zeta_K$, with
 coefficients $\Lambda_K(n)$. Define the Euler--Kronecker constant $\gamma_K$ by
 $\zeta_K'/\zeta_K(s)=-(s-1)^{-1}+\gamma_K+O(s-1)$. Then $\delta=1$, $c=-1-\gamma_K$ and
 \begin{equation}\label{ana:eq:RK}
  R_K(x)=\sum_{n\le x}\Lambda_K(n)\Bigl(\frac1n-\frac1x\Bigr)-\log x+1+\gamma_K .
 \end{equation}
 The setting holds at every $\beta\in(0,1)$, and so it does for every positive combination
 $\sum_jw_jR_{E_j}$ of finitely many fields. An eventual bound
 $\sum_jw_jR_{E_j}(x)\le Cx^{\beta-1}$ excludes zeros of every $\zeta_{E_j}$ in $\Re s>\beta$.
\item \emph{An entire product.} Let $F$ be a finite product of positive integer powers of
 $L$-functions of nontrivial finite-order Hecke characters of number fields, regarded as an
 Euler product over rational primes with all local factors, and with real coefficients in
 $D_F=-F'/F=\sum_{p,k}a_F(p^k)\log p\,p^{-ks}$. Then $\delta=0$, $c=D_F(1)$, and
 $R_F(x)=\sum_{p^k\le x}a_F(p^k)\log p\,(p^{-k}-x^{-1})-D_F(1)$ has no logarithmic term.
\item \emph{An Artin $L$-function with a simple pole.} Let $W$ be a character of a finite Galois
 group for which $L(s,W)$ is known to be holomorphic except for a simple pole at $s=1$, for
 example a permutation character, or $\triv$ plus a nonnegative integer combination of
 characters induced from nontrivial linear characters of subgroups. With all Euler factors,
 $a(p^k)=\log p\cdot\Tr\bigl(g_p^k\mid V_W^{I_p}\bigr)$, where $I_p$ is inertia and $g_p$ a
 Frobenius lift. Here $\delta=1$.
\item \emph{Quotients of zeta functions.} $F_j=Q_j$ is a product of Dedekind zeta functions with
 integer exponents of sum zero that is holomorphic in $\Re s>0$. Here $\delta=0$.
\end{enumerate}
\end{example}

\begin{example}[The source remainder]\label{ana:ex:source}
For a modulus $q\ge1$ and an odd prime $d\nmid q$ write $\chi_d(n)=(n/d)$ and
\[
 \lambda_{d,q}(n)=\anaone_{n\equiv1\ (\mathrm{mod}\ q)}\,\Lambda(n)\bigl(1-\chi_d(n)\bigr)\ge0,\qquad
 R_{d,q}(x)=\sum_{n\le x}\lambda_{d,q}(n)\Bigl(\frac1n-\frac1x\Bigr)
 -\frac{\log x}{\varphi(q)}-c_{d,q}.
\]
This is the remainder bounded from above in the source (its \code{primeError}), and it is the
centred remainder of $\lambda_{d,q}$ in the sense of \cref{ana:def:ramp}. Indeed, let
$E_\chi=L(s,\chi)$ for a nonprincipal character $\chi$ and let $E_\chi$ be the entire function
$(s-1)L(s,\chi)$ for a principal one. Orthogonality and the Euler products give, for $\Re s>1$,
\[
 \sum_{n\ge1}\frac{\lambda_{d,q}(n)}{n^s}
 =\frac1{\varphi(q)}\Bigl\{\frac1{s-1}+\sum_{\chi\bmod q}
 \Bigl(\frac{E'_{\chi\chi_d}(s)}{E_{\chi\chi_d}(s)}-\frac{E'_\chi(s)}{E_\chi(s)}\Bigr)\Bigr\},
\]
so that $\delta=1/\varphi(q)$, and $c_{d,q}$ is real. By \cref{ana:lem:mellin} the transform
$\Phi_{d,q}$ of $R_{d,q}$ is meromorphic on $\Re s>0$ and regular at $1$, its poles are at most
simple, and for $\rho\ne1$
\begin{equation}\label{ana:eq:srcres}
 \Res_{s=\rho}\Phi_{d,q}(s)
 =\frac{\sum_{\chi\bmod q}\ord_\rho E_\chi-\sum_{\chi\bmod q}\ord_\rho E_{\chi\chi_d}}
 {\varphi(q)\,\rho(1-\rho)} .
\end{equation}
Even a multiple zero gives a simple pole, with the multiplicity in the residue. The numerator
is a difference, so $\lambda_{d,q}$ is not in a one-sided setting: an untwisted and a twisted
function vanishing at the same point may cancel. The source first chooses $d$ so that no
twisted function vanishes in a prescribed box. This step is examined in \cref{ana:sec:sparse}.
\end{example}

\begin{proposition}[Converse; \stU]\label{ana:prop:converse}
Let $a$ be in a one-sided setting at $\beta$. Suppose in addition that every $F_j$ is a finite
product of integer powers of $L$-functions of finite-order Hecke characters of number fields,
Dedekind zeta functions included, and that $H$ is bounded on $\Re s\ge\beta+\eps$ for every
$\eps>0$. If no $F_j$ has a zero in $\Re s>\beta$, then for every $\eps>0$
\[
 R(x)=O_\eps\bigl(x^{\beta-1+\eps}\bigr).
\]
Consequently, for $1/2\le\beta<1$ the absence of zeros of all $F_j$ in $\Re s>\beta$ is
equivalent to the family of one-sided bounds $R(x)\le C_\eps x^{\beta-1+\eps}$, $\eps>0$.
\end{proposition}

\begin{proof}
Let $0<\eps<1-\beta$ and $\sigma_0=\beta+\eps$. Each constituent $L$-function, completed and
multiplied by $s(s-1)$ when its character is trivial, is entire of order one. It has a Hadamard
product, and it has $O(\log(\abs t+3))$ zeros with ordinate in $[t-1,t+1]$
\cite[Ch.~5]{IwaniecKowalski2004Analytic}. For $s=\sigma+it$ with $\sigma_0\le\sigma\le2$,
subtract the Hadamard logarithmic derivative at $2+it$ from the one at $s$. Each zero then
contributes a difference of two reciprocals. The zeros in the $k$th unit shell contribute
$O(\log(\abs t+k+3)/k^2)$ in total, the gamma factors obey Stirling's bound, and the Euler series
at $2+it$ converges absolutely. In the signed sum that forms $F_j'/F_j$, a zero $\rho$ with
$\Re\rho>\beta$ has total multiplicity $\ord_\rho F_j=0$ by hypothesis. Every remaining zero
has $\Re\rho\le\beta$, hence distance at least $\eps$ from $s$. Therefore
\[
 A(\sigma+it)=O_\eps\bigl(\log(\abs t+3)\bigr)\qquad(\sigma_0\le\sigma\le2,\ \abs t\ge1),
\]
and at small heights the same bound holds away from $s=1$ by compactness. These are estimates
for fixed functions. No uniformity in the fields or in the degree is claimed.

The inverse Mellin kernel of $(y-1)_+$ gives
\[
 T(x)=\frac1{2\pi i}\int_{2-i\infty}^{2+i\infty}A(s)\frac{x^{s-1}}{s(s-1)}\,\dd s .
\]
Shift the contour to $\Re s=\sigma_0$. The horizontal integrals tend to zero by the logarithmic
bound and the decay $\abs s^{-2}$. The only pole crossed is at $s=1$, with residue
$\delta\log x+c$ by the expansion in the proof of \cref{ana:lem:mellin}. The new vertical
integral is $O_\eps(x^{\sigma_0-1})$, since $\log(\abs t+3)/(1+t^2)$ is integrable. The last
assertion combines this with \cref{ana:thm:zerotest,ana:cor:forms}.
\end{proof}

\subsection{Removal of prime powers}\label{ana:sec:primeonly}
{\raggedright\sources{grh_relative_prime_races_2026-11-25; grh_cyclic_generator_transfer_2026-11-25; grh_a6_transfer_2026-11-25; grh_sylow_transfer_2026-11-25; psl2_grh_transfer; grh_deficit_transfer_2026-11-25}}

The coefficients of a logarithmic derivative at prime powers $p^k$, $k\ge2$, and at finitely
many bad primes, do not affect a test at any exponent $\beta\ge1/2$.

\begin{lemma}[Prime powers and finitely many primes]\label{ana:lem:powers}
Let $a(n)$ be supported on prime powers with $\abs{a(p^k)}\le C_0\log p$, and let $S$ be a finite
set of primes. Call $n=p^k$ \emph{omitted} if $k\ge2$ or $p\in S$. Let $a^\flat(n)=a(n)$ if $n$
is a prime outside $S$ and $a^\flat(n)=0$ otherwise, with series $A^\flat$ and ramp $T^\flat$.
\begin{enumerate}[label=\textup{(\roman*)}]
\item $A-A^\flat=\sum_{n\ \mathrm{omitted}}a(n)n^{-s}$ converges normally on
 $\Re s\ge1/2+\eta$ for every $\eta>0$. Hence $A^\flat$ continues to $\Re s>1/2$ wherever $A$
 does, with the same poles and principal parts.
\item $B=\sum_{n\ \mathrm{omitted}}a(n)/n$ converges absolutely, and
 \begin{equation}\label{ana:eq:powers}
  T(x)-T^\flat(x)=B+O\bigl(x^{-1/2}\bigr).
 \end{equation}
\end{enumerate}
\end{lemma}

\begin{proof}
(i) On $\Re s\ge1/2+\eta$ the terms with $k\ge2$ are dominated by a constant multiple of
$\sum_{n\ge2}\log n/n^{1+2\eta}$, and each of the finitely many local series at $p\in S$
converges normally on $\Re s>0$. (ii) The Chebyshev bound $\sum_{p\le y}\log p=O(y)$ gives
\[
 M(t):=\sum_{p^k\le t,\ k\ge2}\log p=O\bigl(\sqrt t+t^{1/3}\log t\bigr)=O(\sqrt t).
\]
By partial summation,
\[
 \sum_{p^k>x,\ k\ge2}\frac{\abs{a(p^k)}}{p^k}=O\bigl(x^{-1/2}\bigr),\qquad
 \frac1x\sum_{p^k\le x,\ k\ge2}\abs{a(p^k)}=O\bigl(x^{-1/2}\bigr).
\]
A prime $p\in S$ contributes a tail $O(1/x)$ and a correction $O(\log x/x)$.
\end{proof}

For the main statement let $E/\Q$ be a finite Galois extension with group $G$, let $S$ be a
finite set of rational primes containing those ramified in $E$, and let $g_p$ denote the
Frobenius class of $p\notin S$. For a real class function $w$ on $G$ put
$a_w=\langle w,\triv\rangle$ and
\begin{align*}
 D_w(s)&=\sum_{p\notin S}\sum_{k\ge1}w(g_p^k)\log p\;p^{-ks}\quad(\Re s>1),\\
 \Theta_w(x)&=\sum_{p\le x,\ p\notin S}w(g_p)\log p,\qquad
 P_w(x)=\sum_{p\le x,\ p\notin S}w(g_p)\log p\Bigl(\frac1p-\frac1x\Bigr).
\end{align*}
The Chebotarev density theorem with its classical error term gives
\[
 \Theta_w(x)=a_wx+O_E\bigl(x\exp(-c\sqrt{\log x})\bigr)
\]
for some $c>0$; a possible exceptional real zero of the fixed field is absorbed in the error
\cite{LagariasOdlyzko1977Effective}. Since $P_w(x)=\int_1^x\Theta_w(t)t^{-2}\,\dd t$, the limit
\[
 C_w=\lim_{x\to\infty}\bigl(P_w(x)-a_w\log x\bigr)
 =\int_1^\infty\frac{\Theta_w(t)-a_wt}{t^2}\,\dd t
\]
exists, and the \emph{prime-only remainder}
\begin{equation}\label{ana:eq:Rw}
 R_w(x)=P_w(x)-a_w\log x-C_w=-\int_x^\infty\frac{\Theta_w(t)-a_wt}{t^2}\,\dd t
\end{equation}
is continuous, bounded, and tends to zero. This defines the constant without any Riemann
hypothesis. It is a definition, not a finite numerical evaluation.

\begin{theorem}[Prime-only criterion; \stU]\label{ana:thm:primeonly}
With the notation above, suppose that for $\Re s>1$
\[
 D_w(s)=-\sum_{j=1}^rw_j\frac{F_j'(s)}{F_j(s)}+H(s),
\]
where the $w_j$ are positive, every $F_j$ is a finite product of integer powers of
$L$-functions of finite-order Hecke characters of number fields that has no pole in
$1/2<\Re s<1$, and $H$ is holomorphic on $\Re s>1/2$ and bounded on $\Re s\ge1/2+\eta$ for
every $\eta>0$. Let $R$ be the centred remainder of the full coefficient sequence of $D_w$
(\cref{ana:def:ramp}). Then $R(x)=R_w(x)+O(x^{-1/2})$, and the following are equivalent.
\begin{enumerate}[label=\textup{(\roman*)}]
\item No $F_j$ has a zero in $\Re s>1/2$.
\item For every $0<\eps<1/2$ there is $C_\eps$ with $R_w(x)\le C_\eps x^{-1/2+\eps}$ for all
 sufficiently large real $x$.
\item $R_w(x)=O_\eps(x^{-1/2+\eps})$ for every $\eps>0$.
\item For every $0<\eps<1/2$ there is $K_\eps$ such that for all sufficiently large integers
 $n$,
 \begin{equation}\label{ana:eq:dyadicint}
  P_w(n)-P_w(2n)+a_w\log 2\le K_\eps n^{-1/2+\eps}.
 \end{equation}
\item For every $0<\eps<1/2$ there is $K_\eps$ such that for all sufficiently large integers
 $n$,
 \begin{equation}\label{ana:eq:squareint}
  P_w(n)-P_w(n^2)+a_w\log n\le K_\eps n^{-1/2+\eps}.
 \end{equation}
\end{enumerate}
More precisely, for a single $\beta\in(1/2,1)$ an eventual bound with exponent $\beta-1$ in
\textup{(ii)}, \textup{(iv)} or \textup{(v)} excludes zeros of every $F_j$ in $\Re s>\beta$. If
moreover all nontrivial zeros of $\zeta_E$ have real part $1/2$, then
$R_w(x)=O_{E,w}(x^{-1/2}\log^2x)$.
\end{theorem}

\begin{proof}
\emph{Comparison with the full series.} Let $T$ be the ramp of the coefficients
$a(p^k)=w(g_p^k)\log p$, $p\notin S$, of $D_w$. \Cref{ana:lem:powers} gives
\[
 T(x)-P_w(x)=B+O\bigl(x^{-1/2}\bigr),\qquad
 B=\sum_{p\notin S,\ k\ge2}\frac{w(g_p^k)\log p}{p^k}.
\]
The $L$-functions of finite-order Hecke characters are holomorphic and nonvanishing on $\Re s\ge1$
apart from the pole of a zeta function at $1$, so $D_w$ is in a one-sided setting at every
$\beta\in(1/2,1)$. Its residue at $1$ is $\lim_{\sigma\downarrow1}(\sigma-1)D_w(\sigma)=a_w$ by
the asymptotic for $\Theta_w$. As $T(x)-a_w\log x\to C_w+B$, \cref{ana:lem:mellin}(iii)
identifies the centring constant of $T$ as $C_w+B$, and the centred remainder $R$ of $T$
satisfies
\begin{equation}\label{ana:eq:RwR}
 R(x)=R_w(x)+O\bigl(x^{-1/2}\bigr).
\end{equation}

\emph{From a one-sided bound to zeros.} If $R_w(x)\le Cx^{\beta-1}$ eventually, with
$1/2<\beta<1$, then $R(x)\le C'x^{\beta-1}$ eventually by \eqref{ana:eq:RwR}, and
\cref{ana:thm:zerotest} excludes zeros of every $F_j$ in $\Re s>\beta$. Letting $\beta$
decrease to $1/2$ proves (ii)$\Rightarrow$(i).

\emph{From zeros to the two-sided bound.} Under (i), \cref{ana:prop:converse} at $\beta=1/2$ and
\eqref{ana:eq:RwR} give (iii). Clearly (iii) implies (ii), and it implies (iv) and (v), whose
left sides are $R_w(n)-R_w(2n)$ and $R_w(n)-R_w(n^2)$.

\emph{Telescoping.} Assume \eqref{ana:eq:dyadicint} with exponent $\beta-1$, where
$\beta=1/2+\eps$, for all integers $n\ge n_0$. Since $R_w(2^kn)\to0$ unconditionally,
\[
 R_w(n)=\lim_{k\to\infty}\sum_{j=0}^{k-1}\bigl(R_w(2^jn)-R_w(2^{j+1}n)\bigr)
 \le K_\eps n^{\beta-1}\sum_{j\ge0}2^{j(\beta-1)}=\frac{K_\eps}{1-2^{\beta-1}}\,n^{\beta-1}.
\]
Assume instead \eqref{ana:eq:squareint}. Telescoping along the integers $n^{2^i}$ gives
\[
 R_w(n)\le K_\eps\sum_{i\ge0}n^{(\beta-1)2^i}\le2K_\eps n^{\beta-1}
\]
for large $n$, since $2^i\ge i+1$. In both cases only upper bounds are summed, and no unknown
constant has to be computed. Between consecutive integers $R_w$ is differentiable off the
primes, with derivative $\Theta_w(x)x^{-2}-a_wx^{-1}=O_w(x^{-1})$ by the Chebyshev bound. So the
integer bound extends to all sufficiently large real $x$ with a changed constant, which is (ii)
with the same $\beta$.

\emph{The sharper bound.} If $\zeta_E$ satisfies the Riemann hypothesis, the effective
Chebotarev theorem gives $\Theta_w(x)=a_wx+O_{E,w}(\sqrt x\log^2x)$
\cite{LagariasOdlyzko1977Effective,GrenieMolteni2019Explicit}, and \eqref{ana:eq:Rw} gives the
bound for $R_w$.
\end{proof}

\begin{remark}[The statistics in \textup{(iv)} and \textup{(v)} are finite]\label{ana:rem:finite}
The constant $C_w$ does not enter \eqref{ana:eq:dyadicint} or \eqref{ana:eq:squareint}. Direct
subtraction gives
\[
 P_w(n)-P_w(2n)=-\frac{\Theta_w(n)}{2n}
 -\sum_{\substack{n<p\le2n\\p\notin S}}w(g_p)\log p\Bigl(\frac1p-\frac1{2n}\Bigr).
\]
A proof of \eqref{ana:eq:dyadicint} need not evaluate any $L'(1)/L(1)$. It has to control a
signed combination of a cumulative count and a short-interval count. Replacing either
inequality by its negative does not give the same conclusion, because the sign of the real
residues would be reversed (\cref{ana:rem:sign}).
\end{remark}

\begin{remark}[Scope]\label{ana:rem:base}
\Cref{ana:thm:primeonly} is the common analytic part of the criteria of
\cref{sec:nab,sec:grp}. The choice of $w$ and the verification of the displayed identity for
$D_w$ belong to those sections. When the nonvanishing of the $F_j$ is known to be equivalent to
the Riemann hypothesis for $\zeta_E$, each of (ii)--(v) is equivalent to that hypothesis. The
excluded set $S$ may be any finite set containing the ramified primes. When $w$ is given by the
number of roots modulo $p$ of a polynomial, $S$ may consist of the primes dividing its
discriminant, and (ii)--(v) refer to rational primes and root counts only. The same statements
hold over a base number field $k$ in place of $\Q$, with prime ideals of $k$ ordered by absolute
norm: the Mellin calculation depends only on the norms, and each rational prime lies below at
most $[k:\Q]$ prime ideals, so \cref{ana:lem:powers} is unchanged.
\end{remark}

\subsection{A finite-difference detector}\label{ana:sec:fd}
{\raggedright\sources{grh_deficit_transfer_2026-11-25; grh_a6_transfer_2026-11-25}}

The integer-scale statistics of \cref{ana:thm:primeonly} avoid the centring constant by
telescoping, which uses the limit $R_w(x)\to0$. A difference at two scales of fixed ratio avoids
the constant without any limit, and accepts hypotheses in mean. Let $a$ be as in
\cref{ana:def:ramp}, fix $\lambda>1$, and put
\begin{equation}\label{ana:eq:fd}
 \Delta_\lambda(x)=T(x)-T(\lambda x)+\delta\log\lambda=R(x)-R(\lambda x)\qquad(x\ge1).
\end{equation}
Term by term, including at $n=x$ and $n=\lambda x$,
\[
 T(\lambda x)-T(x)=\frac1x\sum_{n\le\lambda x}a(n)K_\lambda\Bigl(\frac nx\Bigr),\qquad
 K_\lambda(y)=\begin{cases}1-\lambda^{-1},&0<y\le1,\\ y^{-1}-\lambda^{-1},&1<y\le\lambda,\\
 0,&y>\lambda,\end{cases}
\]
with a nonnegative kernel. The statistic uses the coefficients up to $\lambda x$ only, and no
value of a logarithmic derivative at $1$.

\begin{lemma}[Mellin identity]\label{ana:lem:fd}
Suppose that $A$ extends meromorphically to $\Re s>\beta$, where $0<\beta<1$, with at most a
simple pole at $1$, of residue $\delta$. For $\Re s>1$,
\begin{gather*}
 \int_1^\infty\Delta_\lambda(x)x^{-s}\,\dd x
 =-H_\lambda(s)A(s)+\frac{\delta\log\lambda}{s-1}+J_\lambda(s),\\
 H_\lambda(s)=\frac{\lambda^{s-1}-1}{s(s-1)},\qquad
 J_\lambda(s)=\lambda^{s-1}\int_1^\lambda T(u)u^{-s}\,\dd u .
\end{gather*}
The function $J_\lambda$ is entire and the right side is regular at $s=1$, where
$H_\lambda(1)=\log\lambda$. The function $H_\lambda$ has no zero in $0<\Re s<1$, and
$H_\lambda(\sigma)>0$ for real $\sigma>0$.
\end{lemma}

\begin{proof}
By \eqref{ana:eq:mellin}, $\int_1^\infty T(x)x^{-s}\,\dd x=A(s)/(s(s-1))$. In the term with
$T(\lambda x)$ substitute $u=\lambda x$; the omitted interval $[1,\lambda]$ gives $J_\lambda$.
Since $\lambda^{s-1}=1$ forces $\Re s=1$, the function $H_\lambda$ has no zero in the open
strip, and $\lambda^{\sigma-1}-1$ has the sign of $\sigma-1$.
\end{proof}

The kernel $1/(s(s-1))$ of the undifferenced ramp is negative on $(0,1)$, and so is
$-H_\lambda$. Thus $R$ and $\Delta_\lambda$ detect poles of $A$ with the same sign of residue,
while the reversed difference $T(\lambda x)-T(x)$ detects the opposite sign.

\begin{theorem}[One-sided averaged detector; \stU]\label{ana:thm:fd}
In the situation of \cref{ana:lem:fd}, assume that all poles of $A$ in $\Re s>\beta$ other than
$1$ are simple, lie in $\beta<\Re s<1$, and have nonpositive residue whenever they are real. If
\begin{equation}\label{ana:eq:fdint}
 \int_1^\infty\anapos{\Delta_\lambda(x)}x^{-\sigma}\,\dd x<\infty
 \qquad\text{for every real }\sigma>\beta,
\end{equation}
then $A$ has no pole in $\Re s>\beta$ except possibly at $1$. It suffices that
$\int_X^{2X}\anapos{\Delta_\lambda(x)}\,\dd x=O_\eps(X^{\beta+\eps})$ for every $\eps>0$. The
pointwise condition $\Delta_\lambda(x)\le C_\eps x^{\beta-1+\eps}$ is stronger and also suffices.
\end{theorem}

\begin{proof}
The function $\Delta_\lambda$ is continuous and $O(\log 2x)$. By \cref{ana:lem:fd} its
transform is meromorphic on $\Re s>\beta$ and regular on $\Re s\ge1$. Its poles in
$\beta<\Re s<1$ are those of $A$, with residue $-H_\lambda(\rho)\Res_{s=\rho}A$, which is
nonnegative at real $\rho$. \Cref{ana:thm:onesided} shows that the transform is holomorphic on
$\Re s>\beta$. So is $H_\lambda A$ on $\beta<\Re s<1$, and $H_\lambda$ does not vanish there.
The two sufficient conditions are those of \cref{ana:cor:forms}.
\end{proof}

In a one-sided setting (\cref{ana:def:setting}) the hypothesis on the poles holds, so
\cref{ana:thm:zerotest} remains true with $R(x)-R(\lambda x)$ in place of $R(x)$, and the
constant $c$ is never needed. Under the hypotheses of \cref{ana:prop:converse} one has
conversely $\Delta_\lambda(x)=O_\eps(x^{\beta-1+\eps})$, and by \cref{ana:lem:powers} the
passage to primes changes $\Delta_\lambda$ by $O(x^{-1/2})$. For weights with the opposite
orientation (\cref{ana:rem:sign}) the theorem applies to $-a$: the detector is then an upper
bound for $T(\lambda x)-T(x)-\delta\log\lambda$. Two separate upper bounds at $x$ and at
$\lambda x$ cannot be subtracted to bound $\Delta_\lambda$. An arithmetic construction has to
produce the increment itself, or a matching lower bound for the subtracted term.

\subsection{Signed averages and joint filters}\label{ana:sec:mean}
{\raggedright\sources{grh_method_transfers_2026-11-25_2; grh_transfer_synthesis_2026-11-25}}

By \cref{ana:cor:forms} a bound for the positive part in mean suffices. This subsection
determines how far cancellation may be admitted. Throughout, $R$ and $F$ satisfy the analytic
hypotheses of \cref{ana:thm:onesided}.

\begin{theorem}[Signed moving window; \stU]\label{ana:thm:moving}
Fix $h>1$. If, for some $C\ge0$,
\begin{equation}\label{ana:eq:moving}
 \int_x^{hx}R(t)\,\dd t\le Cx^{\beta}\qquad\text{for every sufficiently large real }x,
\end{equation}
then $F$ is holomorphic on $\Re s>\beta$. It also suffices to prove \eqref{ana:eq:moving} with
exponent $\beta+\eps$ for every $\eps>0$, with constants depending on $\eps$.
\end{theorem}

\begin{proof}
Set $T_hR(x)=x^{-1}\int_x^{hx}R(t)\,\dd t$. Fubini's theorem, far to the right, gives
\[
 \int_1^\infty T_hR(x)x^{-s}\,\dd x=\frac{h^s-1}{s}F(s)
 -\frac1s\int_1^hR(t)\bigl(h^st^{-s}-1\bigr)\,\dd t ,
\]
because the inner integral in the reversed order is $\int_{\max(1,t/h)}^tx^{-s-1}\,\dd x$. The
finite-interval correction is entire. The multiplier $(h^s-1)/s$ is entire, has no zero on
$\Re s>0$, and is positive for real $s>0$. Thus the transform of $T_hR$ has exactly the poles of
$F$ in $\Re s>\beta$, and its real residues are positive.

By \eqref{ana:eq:moving}, $T_hR(x)\le Cx^{\beta-1}$ eventually, and $T_hR$ is continuous.
\Cref{ana:thm:onesided} with \cref{ana:cor:forms}(a) shows that the transform of $T_hR$ has no
pole in $\Re s>\beta$, and neither has $F$. The last claim follows by applying this at each
$\beta+\eps$.
\end{proof}

\begin{remark}\label{ana:rem:box}
(1) For real $\nu$ the weighted box $T_{h,\nu}R(x)=x^{-1-\nu}\int_x^{hx}t^\nu R(t)\,\dd t$ has
multiplier $(h^{s+\nu}-1)/(s+\nu)$, whose zeros lie on $\Re s=-\nu$. The criterion remains valid
on $\Re s>\beta$ whenever $\beta+\nu\ge0$. For an arbitrary varying weight the transform must be
computed afresh.

(2) For a centred ramp $R$ one has $\abs{R(x)}=O(\log^2(2x))$, so changing $x$ by at most one
changes the integral in \eqref{ana:eq:moving} by $O(\log^2(2x))$, and integer starting points
suffice. Exponentially sparse starting points do not, as the next example shows.
\end{remark}

\begin{example}[Signed averages on the dyadic grid are insufficient]\label{ana:ex:dyadic}
Fix $0<\beta<a<1$, put $\tau=2\pi/\log2$ and
\[
 R(x)=\frac{\dd}{\dd x}\bigl(x^a\sin(\tau\log x)\bigr)
 =x^{a-1}\bigl(a\sin(\tau\log x)+\tau\cos(\tau\log x)\bigr).
\]
Then $\int_{2^k}^{2^{k+1}}R(x)\,\dd x=0$ for every integer $k\ge0$, while
\[
 F(s)=\frac{\tau s}{(s-a)^2+\tau^2}\qquad(\Re s>a)
\]
has simple poles at $a\pm i\tau$, to the right of $\beta$, and no real pole. All analytic
hypotheses of \cref{ana:thm:onesided} hold. Here
$\int_{2^k}^{2^{k+1}}R(x)_+\,\dd x=C_a2^{ka}$ with $C_a>0$, by the substitution $x=2^ke^u$. So
cancellation in signed averages over a fixed partition can hide exactly the poles that the test
is meant to exclude.
\end{example}

A positive averaging kernel is not enough either: its Mellin multiplier must have no zero in
the half-plane. The following theorem makes this exact for averages over geometric scales.
Fix $h>1$. For a nonzero polynomial $P(z)=\sum_{k=0}^dp_kz^k$ with real coefficients $p_k\ge0$,
put
\[
 T_PR(x)=\sum_{k=0}^dp_kR(h^kx),\qquad z(s)=h^{s-1}.
\]
A change of variables in each term gives, far to the right,
\begin{equation}\label{ana:eq:filter}
 \int_1^\infty T_PR(x)x^{-s}\,\dd x=P(z(s))F(s)-E_P(s),\qquad
 E_P(s)=\sum_kp_kh^{k(s-1)}\int_1^{h^k}R(u)u^{-s}\,\dd u,
\end{equation}
with $E_P$ entire. At a real point $s$ one has $P(h^{s-1})>0$.

\begin{theorem}[Joint detection; \stU]\label{ana:thm:joint}
Let $P_1,\dots,P_r$ be nonzero polynomials with nonnegative real coefficients, and suppose that
the functions $P_j(z(s))$ have no common zero in $\Re s>\beta$. If, for each $j$,
\begin{equation}\label{ana:eq:joint}
 \int_1^\infty\anapos{T_{P_j}R(x)}x^{-\sigma}\,\dd x<\infty\qquad(\sigma>\beta),
\end{equation}
then $F$ is holomorphic on $\Re s>\beta$. Conversely, the absence of a common zero is necessary
for this implication to hold for every $R$ satisfying the analytic hypotheses of
\cref{ana:thm:onesided}.
\end{theorem}

\begin{proof}
By \eqref{ana:eq:filter} the transform of $T_{P_j}R$ has at most simple real poles with positive
residues, because its multiplier is positive at real points. \Cref{ana:thm:onesided} applies to
each filtered function, so every $P_j(z(s))F(s)$ is holomorphic on $\Re s>\beta$. At any point
of the half-plane some multiplier is nonzero, and division by it shows that $F$ is holomorphic
there.

For the necessity let $\rho=a+i\tau$ be a common zero with $a>\beta$. It is not real. The
function
\[
 R(x)=\Re x^{\rho-1}=x^{a-1}\cos(\tau\log x)
\]
has a meromorphic Mellin transform with simple poles at $\rho,\bar\rho$ and no real pole, yet
$T_{P_j}R(x)=\Re\bigl(P_j(h^{\rho-1})x^{\rho-1}\bigr)=0$ for every $j$. All hypotheses
\eqref{ana:eq:joint} hold and $F$ is not holomorphic on $\Re s>\beta$.
\end{proof}

The necessity concerns the whole analytic class. It does not assert that a given common zero is
a zero of a particular $L$-function. The simplest instance is the single filter
$P(z)=1+2^{1-a}z$ with $h=2$: for $R(x)=x^{a-1}\cos(\pi\log x/\log2)$ the positive combination
$R(x)+2^{1-a}R(2x)$ vanishes identically. If a single available multiplier is zero-free on
$\Re s>\beta$, as in \cref{ana:thm:moving,ana:thm:fd}, it suffices on its own.

\begin{example}[Two filters with integer weights]\label{ana:ex:filters}
Take $h=2$, $P(z)=5+6z$ and $Q(z)=25+30z+36z^2$. The identity
\begin{equation}\label{ana:eq:bezout}
 Q(z)-6zP(z)=25
\end{equation}
shows that $P$ and $Q$ have no common zero. Hence the statistics
\[
 A_R(x)=5R(x)+6R(2x),\qquad B_R(x)=25R(x)+30R(2x)+36R(4x)
\]
suffice jointly: if, for every $\eps>0$ and all large dyadic $X$,
\[
 \int_X^{2X}\anapos{A_R(x)}\,\dd x\ll_\eps X^{\beta+\eps},\qquad
 \int_X^{2X}\anapos{B_R(x)}\,\dd x\ll_\eps X^{\beta+\eps},
\]
then $F$ is holomorphic on $\Re s>\beta$. Moreover the exact identity
$25R(x)=B_R(x)-6A_R(2x)$, with the absolute integrability of both filtered functions supplied
by \cref{ana:thm:onesided}, shows that the Mellin integral of $R$ converges absolutely there. At
$\beta=1/2$ each displayed bound may be replaced by a square moment $\ll_\eps X^{\eps}$ of the
positive part. Each filter alone has zeros on the line
$\Re s=1+\log(5/6)/\log2\approx0.7369656$: for $P$ at the ordinates $(2k+1)\pi/\log2$, for $Q$
at $(2k\pm2/3)\pi/\log2$. So each can conceal poles to the right of $1/2$, and the certificate
\eqref{ana:eq:bezout}, not a numerical root search, shows that the pair cannot. Integer
weights are chosen because direct sums and repeated copies of integral lattices produce
integer-weighted sums of logarithmic covolumes.
\end{example}

\begin{proposition}[Positive coefficients do not repair a single filter; \stU]\label{ana:prop:onefilter}
There are positive, uniformly bounded real coefficients $a_n$ and constants $A_0>0$, $c$ such
that the remainder $R(x)=\sum_{n\le x}a_n(1/n-1/x)-A_0\log x-c$ satisfies $5R(x)+6R(2x)\le0$
for all $x\ge1$, while its Mellin transform, which is meromorphic on $\Re s>0$ with no real pole
there, has poles with real part $1+\log(5/6)/\log2>1/2$.
\end{proposition}

\begin{proof}
Let $a=1+\log(5/6)/\log2$, $\tau=\pi/\log2$, $z=a-1+i\tau$ and $R_0(x)=\epsilon\Re x^z$. Given
$A_0>0$ choose $\epsilon>0$ with $\epsilon\abs{z(z+1)}\le A_0/2$ and
$\epsilon\abs{\Re z}\le A_0/2$. The measure
\[
 \mu=(A_0+\epsilon\Re z)\,\delta_1
 +\bigl[A_0+\epsilon\Re\bigl(z(z+1)x^z\bigr)\bigr]\dd x
\]
is positive with bounded density, and its mass below $x$ is $A_0x+\epsilon\Re(zx^{z+1})$.
Differentiating the ramp $S_\mu(x)=\int_{[1,x]}(u^{-1}-x^{-1})\,\dd\mu(u)$ and matching the
value at $x=1$ gives $S_\mu(x)-A_0\log x+\epsilon=R_0(x)$. The filter $5+6z$ annihilates $R_0$,
since $2^{z}=-5/6$, although the transform of $R_0$ has poles at $a\pm i\tau$. Now move the
density on each interval $[n,n+1)$ to the integer $n$ and correct the limiting constant, as in
the proof of \cref{ana:prop:mptsharp}. The computation \eqref{ana:eq:discreteid} there shows that
the remainder becomes $R=R_0-D$ with $0\le D(x)\ll1/x$. Then $5R(x)+6R(2x)=-5D(x)-6D(2x)\le0$,
and the poles persist because the transform of $D$ is holomorphic on $\Re s>0$.
\end{proof}

The comparison of covolumes in the source (\cref{src:sec:mechanism}) produces its bound for
the remainder at each $x$ separately. The next statement records how its losses may be
averaged.

\begin{proposition}[Rank-sensitive loss in mean; \stU]\label{ana:prop:rankloss}
At a given $x$ let $d>0$, $0<n\le d$, $M>0$, $H\ge0$ and $L,U,I\ge0$, and suppose that a real
number $V$ satisfies
\[
 d(H-L)\le V+2dI,\qquad V\le n\bigl(H-MR(x)+U\bigr).
\]
With $\alpha=n/d$,
\begin{equation}\label{ana:eq:rank}
 R(x)_+\le\frac{\bigl(L+\alpha U+2I-(1-\alpha)H\bigr)_+}{\alpha M}
 \le\frac1M\Bigl(U+\frac{L+2I}{\alpha}\Bigr).
\end{equation}
Consequently, if these data exist for all large $x$, with $M(x)\ge c_0x^{\theta}$ for some
$0<\theta<1$, and if for every $\eps>0$
\begin{equation}\label{ana:eq:rankmean}
 \int_X^{2X}\Bigl(U(x)+\frac{L(x)+2I(x)}{\alpha(x)}\Bigr)\dd x\ll_\eps X^{1+\eps},
\end{equation}
then $\int_X^{2X}R(x)_+\,\dd x\ll_\eps X^{1-\theta+\eps}$, and under the analytic hypotheses of
\cref{ana:thm:onesided} the transform of $R$ has no pole in $\Re s>1-\theta$.
\end{proposition}

\begin{proof}
Combining the two inequalities and dividing by $d$ gives
$\alpha MR(x)\le L+\alpha U+2I-(1-\alpha)H$. Take positive parts and divide by $\alpha M$. As
$H\ge0$ and $\alpha\le1$, dropping the nonpositive term gives the second inequality. The rest
is \cref{ana:cor:forms}(b).
\end{proof}

No uniform lower bound for the spanning fraction $\alpha$ is required. For conjugate exponents
$p,q>1$, H\"older's inequality derives \eqref{ana:eq:rankmean} from
$\int_X^{2X}U\ll_\eps X^{1+\eps}$, $\int_X^{2X}(L+2I)^p\ll_\eps X^{1+\eps}$ and
$\int_X^{2X}\alpha^{-q}\ll_\eps X^{1+\eps}$. The bounds must be combined at each $x$ before
averaging. Separate bounds for the averages of $L$, $I$ and $\alpha$ are insufficient: on a
probability space let $L=1$, and let $\alpha=\eta^2$ with probability $\eta$ and $\alpha=1$
otherwise. Then $\anaEx L=1$ and $\anaEx\alpha\to1$ as $\eta\to0$, but
$\anaEx(L/\alpha)=1-\eta+\eta^{-1}\to\infty$.

\subsection{What averaged bounds recover pointwise}\label{ana:sec:mpt}
{\raggedright\sources{grh_transfer_synthesis_2026-11-25; grh_method_transfers_2026-11-25_2}}

A bound in mean for the positive part is a weaker input than a pointwise bound. For
remainders of positive measures one can say exactly how much weaker. Estimates from $L^p$ to
$L^\infty$ for semiconvex functions are in the literature
\cite[Proposition~8.2]{BragaFigalliMoreira2019Optimal}; the one-sided statement, its constants
and the countermodel below are proved directly. Let $\mu$ be a positive
locally finite Borel measure on $[1,\infty)$, let $A_0\ge0$ and $c\in\R$, and put
\begin{equation}\label{ana:eq:posramp}
 S(x)=\int_{[1,x]}\Bigl(\frac1u-\frac1x\Bigr)\dd\mu(u),\qquad
 R(x)=S(x)-A_0\log x-c\qquad(x\ge1).
\end{equation}
The centred ramp of nonnegative coefficients $a(n)$ has this form with $A_0=\delta$, and the
source remainder has it with $A_0=1/\varphi(q)$ and its nonnegative prime-power weights. For
$p\ge1$ and $x\ge3/2$ define
\[
 M_p(x)=\frac1x\int_{2x/3}^{2x}R(y)_+^p\,\dd y,\qquad h=R(x)_+,\qquad c_p=\frac{2^{2-p}}9 .
\]

\begin{theorem}[Positive mass forced by a positive value; \stU]\label{ana:thm:mpt}
If $A_0>0$, then
\begin{equation}\label{ana:eq:mptmass}
 M_p(x)\ge c_ph^p\min\Bigl\{1,\sqrt{\frac{3h}{A_0}}\Bigr\}.
\end{equation}
Consequently, with $D_p=9\cdot2^{p-2}$,
\begin{equation}\label{ana:eq:mptpoint}
 h\le\max\Bigl\{\bigl(D_pM_p(x)\bigr)^{1/p},
 \bigl(D_p\sqrt{A_0/3}\,M_p(x)\bigr)^{1/(p+1/2)}\Bigr\}.
\end{equation}
For $A_0=0$ one has $M_p(x)\ge c_ph^p$.
\end{theorem}

\begin{proof}
The function $t\mapsto\int_{[1,\infty)}(u^{-1}-t)_+\,\dd\mu(u)$ is finite and convex for $t>0$,
and equals $S(1/t)$ for $0<t\le1$. Set $f(v)=R(x/v)$ for $1/2\le v\le3/2$. Then
$f(v)-A_0\log v$ is convex, and its midpoint inequality gives
\[
 f(1-z)+f(1+z)\ge2f(1)+A_0\log(1-z^2)\qquad(0\le z\le1/2).
\]
The assertion is immediate if $h=0$. Otherwise $f(1)=h>0$. Since $t\mapsto t_+^p$ is convex and
nondecreasing,
\[
 f(1-z)_+^p+f(1+z)_+^p\ge2\Bigl(h+\frac{A_0}2\log(1-z^2)\Bigr)_+^p .
\]
As $\log(1-z^2)\ge-4z^2/3$ for $0\le z\le1/2$, the right side is at least $2(h/2)^p$ on
$0\le z\le r$, where $r=\min\{1/2,\sqrt{3h/(4A_0)}\}$. The change of variables $y=x/v$ yields
\[
 M_p(x)=\int_{1/2}^{3/2}f(v)_+^p\frac{\dd v}{v^2}\ge\frac49\int_{1/2}^{3/2}f(v)_+^p\,\dd v
 \ge\frac{8r}9\Bigl(\frac h2\Bigr)^p,
\]
which is \eqref{ana:eq:mptmass}. Splitting at $h=A_0/3$ gives \eqref{ana:eq:mptpoint}. If
$A_0=0$ take $r=1/2$. No differentiability of $R$ or of $\mu$ is used.
\end{proof}

\begin{corollary}[Dyadic moments to pointwise values]\label{ana:cor:mpt}
Suppose $A_0>0$, $p\ge1$ and $\gamma>0$. If for every $\eps>0$ and every sufficiently large
dyadic $X$
\begin{equation}\label{ana:eq:mptdyadic}
 \int_X^{2X}R(y)_+^p\,\dd y\ll_\eps X^{1-p\gamma+\eps},
\end{equation}
then $R(x)_+\ll_\eps x^{-2p\gamma/(2p+1)+\eps}$. At the square-root scale $\gamma=1/2$ the
pointwise exponent is $-p/(2p+1)$: it is $-1/3$ from first moments and $-2/5$ from second
moments.
\end{corollary}

\begin{proof}
Cover $[2x/3,2x]$ by a bounded number of dyadic intervals of size comparable with $x$. Then
$M_p(x)\ll_\eps x^{-p\gamma+\eps}$, so $M_p(x)\to0$, and \eqref{ana:eq:mptmass} forces $h\to0$.
The small-$h$ branch of \eqref{ana:eq:mptpoint} then applies.
\end{proof}

Under the analytic hypotheses of \cref{ana:thm:onesided} each input \eqref{ana:eq:mptdyadic}
with $\gamma=1/2$ already excludes poles to the right of $1/2$, by \cref{ana:cor:forms} and
H\"older's inequality. Passing through a pointwise estimate first would weaken the conclusion
to the half-plane $\Re s>1-p/(2p+1)$. \Cref{ana:thm:mpt} is a local certificate, and the next
result shows that its exponent is the limit of what positivity alone gives.

\begin{proposition}[Optimal exponent in the positive-ramp class; \stU]\label{ana:prop:mptsharp}
Fix $p\ge1$ and $0<\gamma<1$, and set $\alpha=2p\gamma/(2p+1)$. There are positive coefficients
$a_n$, with $1/2\le a_n\le3/2$ for $n\ge2$ and $\sum_{n\le x}a_n\sim x$, and a real constant
$c_d$, such that
\[
 R_d(x)=\sum_{n\le x}a_n\Bigl(\frac1n-\frac1x\Bigr)-\log x-c_d
\]
satisfies \eqref{ana:eq:mptdyadic} with $\eps=0$, while $R_d(4^k)\sim(4^k)^{-\alpha}$ as
$k\to\infty$. Thus the exponent in \cref{ana:cor:mpt} cannot be increased for this class.
\end{proposition}

\begin{proof}
Let $X_k=4^k$, $h_k=X_k^{-\alpha}$ and $L_k=10X_k\sqrt{h_k}$, and start at an integer $k_0$
with $h_k\le1/400$ for $k\ge k_0$. Let $\psi(z)=(1-z^2)^3$ for $\abs z\le1$ and $\psi(z)=0$
otherwise, a function of class $C^2$, and put
\[
 R_c(x)=\sum_{k\ge k_0}h_k\,\psi\Bigl(\frac{x-X_k}{L_k}\Bigr).
\]
The supports are disjoint, because $L_k\le X_k/2$ and successive centres have ratio four.
Since $\psi''=-6+36z^2-30z^4$ on the support, one finds $\max\abs{\psi'}=96/(25\sqrt5)<2$ and
$\max\abs{\psi''}=6$. Put $w(x)=1+(x^2R_c'(x))'$ and $\mu_c=\delta_1+w(x)\,\dd x$ on
$[1,\infty)$. On the $k$th support $x\le3X_k/2$, hence
\[
 \abs{w(x)-1}\le\frac{6x^2h_k}{L_k^2}+\frac{4xh_k}{L_k}
 \le\frac{27}{200}+\frac35\sqrt{h_k}\le\frac{33}{200},
\]
so $1/2\le w\le3/2$. Integration gives $\mu_c([1,x])=x+x^2R_c'(x)$. Differentiating the ramp and
comparing values at $1$ proves the exact identity
\[
 S_{\mu_c}(x)-\log x=R_c(x).
\]
On each bump
\[
 \int R_c(x)^p\,\dd x=h_k^pL_k\int_{-1}^1\psi^p
 =10\Bigl(\int_{-1}^1\psi^p\Bigr)X_k^{1-\alpha(p+1/2)}=C_pX_k^{1-p\gamma},
\]
and every dyadic interval meets at most two supports, with comparable centres. This is the mean
bound for $R_c$, whose value at $X_k$ is exactly $X_k^{-\alpha}$.

To pass to integer support set
\[
 a_n=\anaone_{n=1}+\int_n^{n+1}w(u)\,\dd u,\qquad
 c_d=\int_1^\infty w(u)\Bigl(\frac1{\floor u}-\frac1u\Bigr)\dd u .
\]
The last integral converges absolutely, and $1/2\le a_n\le3/2$ for $n\ge2$. For
$N=\floor x$, direct subtraction gives
\begin{equation}\label{ana:eq:discreteid}
 R_d(x)-R_c(x)=-\int_{N+1}^\infty w(u)\Bigl(\frac1{\floor u}-\frac1u\Bigr)\dd u
 +\int_x^{N+1}w(u)\Bigl(\frac1u-\frac1x\Bigr)\dd u .
\end{equation}
Both terms are nonpositive, and for $x\ge2$ their absolute values sum to at most
$3/(2N)+3/(4x^2)\le4/x$. Hence
\begin{equation}\label{ana:eq:discreteerr}
 0\le R_c(x)-R_d(x)\le4/x ,
\end{equation}
so that $(R_d)_+\le R_c$, which preserves the mean bound. Since $\alpha<1$,
$R_d(X_k)=X_k^{-\alpha}+O(X_k^{-1})\sim X_k^{-\alpha}$. Finally $x^2R_c'(x)=O(x\sqrt{h_k})=o(x)$
on the supports and vanishes elsewhere, so $\mu_c([1,x])\sim x$, and rounding changes this by
$O(1)$.
\end{proof}

The coefficients of \cref{ana:prop:mptsharp} obey a stronger size bound than $O(\log n)$, but
they have no prime-power support, no Euler product and no functional equation. The example
limits an abstract inference. It says nothing against sharper statements about primes. It
contains the earlier examples of narrow positive spikes, which separate mean from pointwise
control without the ramp structure. The model also has the analytic properties required by
the mean criterion.

\begin{proposition}[Compatibility with meromorphic continuation; \stU]\label{ana:prop:mptmero}
In the construction of \cref{ana:prop:mptsharp} write $(s)_j=s(s+1)\cdots(s+j-1)$ and
\[
 J_m(s)=\frac{(s)_{2m}}{(2m)!}\int_{-1}^1z^{2m}\psi(z)\,\dd z .
\]
Then the Mellin transform of $R_c$ is
\begin{equation}\label{ana:eq:mptmero}
 F_c(s)=10\sum_{m\ge0}100^mJ_m(s)\,
 \frac{4^{k_0(1-s-\alpha(m+3/2))}}{1-4^{1-s-\alpha(m+3/2)}} ,
\end{equation}
a meromorphic function on $\C$ whose possible poles are simple and lie at
$s=1-\alpha(m+3/2)+2\pi in/\log4$, $m\ge0$, $n\in\Z$. At a positive real pole the residue is
positive. The rightmost line of poles is
\[
 \Re s=1-\frac{3\alpha}2=1-\frac{3p\gamma}{2p+1}\le1-\gamma ,
\]
with equality when $p=1$. The transform $F_d$ of $R_d$ differs from $F_c$ by a function
holomorphic on $\Re s>0$.
\end{proposition}

\begin{proof}
With $\delta_k=L_k/X_k=10X_k^{-\alpha/2}\le1/2$, expand $(1+\delta_kz)^{-s}$ in the integral
over the $k$th bump and use that $\psi$ is even. This gives \eqref{ana:eq:mptmero} far to the
right. On a compact set avoiding the displayed points, the late terms are bounded by a
polynomial in $m$ times $(100\cdot4^{-k_0\alpha})^m$, and this ratio is at most $1/4$. So the
series is meromorphic on $\C$. At a positive real pole $u=1-\alpha(m+3/2)$ the residue is
$10\cdot100^mJ_m(u)/\log4>0$. The last assertion follows from \eqref{ana:eq:discreteerr}.
Global continuation of $F_d$ is not asserted.
\end{proof}

At $\gamma=1/2$ these models meet the mean budget exactly and have no pole to the right of
$1/2$, yet their peaks decay only like $x^{-p/(2p+1)}$. The sharper pointwise bounds known for
arithmetic $L$-functions under their Riemann hypotheses use additional structure. For a
construction that controls a remainder in mean, the mean criterion should therefore be applied
directly, and \cref{ana:thm:mpt} kept for diagnostics on finite windows.

\subsection{Auxiliary windows}\label{ana:sec:window}
{\raggedright\sources{grh_cyclic_generator_transfer_2026-11-25; grh_kummer_window_transfer_2026-11-25; grh-progression-densities-and-cubic-prime-race; grh_deficit_transfer_2026-11-25}}

The comparison in the source involves the remainder at a second, much larger scale. For two
error functions $E$ and $E'$ of the logarithmic variable, the source selects a point
$u\in[t,t^2]$ with
\[
 E(u)-XE'(u)\le e^{-t},\qquad\text{uniformly for }0\le X\le e^t
\]
(its lemma \code{IsPolar.exists_window_sub_mul_le}; see \cref{src:sec:mechanism}). The results
of this
subsection replace the window $[t,t^2]$ by a linear window $[t,Lt]$, treat finitely many errors
with one averaging measure, allow exponentially large coefficients, and determine which
windows work for the whole analytic class. They use no Riemann hypothesis.

\begin{theorem}[Common truncated Gamma average; \stU]\label{ana:thm:window}
Let $E_1,\dots,E_q\colon[0,\infty)\to\R$ be measurable with $\abs{E_j(v)}\le Be^{bv}$ for
constants $B>0$ and $b\ge0$. Suppose that for some $a>b$ the Laplace transforms
$F_j(z)=\int_0^\infty E_j(v)e^{-zv}\,\dd v$ continue holomorphically to a neighbourhood of the
closed disc $\abs{z-a}\le a$. Then for every $\gamma>0$ there are $L>1$ and $\kappa>0$ such
that, for all sufficiently large $t$, the probability measure
\[
 \dd\nu_t(v)=\frac{\anaone_{[t,Lt]}(v)\,v^{\ceil{\kappa t}}e^{-av}\,\dd v}
 {\int_t^{Lt}u^{\ceil{\kappa t}}e^{-au}\,\dd u}
\]
satisfies $\bigl|\int E_j\,\dd\nu_t\bigr|\le e^{-\gamma t}$ for every $j$. If the $E_j$ are
continuous, there are at most $q+1$ points of $[t,Lt]$ and nonnegative weights of sum one,
independent of $j$, with the same bound for all the weighted averages.
\end{theorem}

\begin{proof}
By compactness and finiteness of the family there is a radius $R>a$ such that all $F_j$ are
holomorphic on a neighbourhood of $\abs{z-a}\le R$. Let $C$ bound their absolute values there.
The Cauchy estimate and differentiation under the integral at $a$ give
\[
 \Bigl|\int_0^\infty v^ke^{-av}E_j(v)\,\dd v\Bigr|=\abs{F_j^{(k)}(a)}\le Ck!\,R^{-k}.
\]
Let $\mu_k$ be the Gamma probability measure $a^{k+1}v^ke^{-av}\,\dd v/k!$. The full mean of
$E_j$ is bounded by $Ca(a/R)^k$. For the low tail, with $e^{-av}$ as a Markov test on $v<t$,
\[
 \int_{v<t}\abs{E_j}\,\dd\mu_k\le Be^{(a+b)t}2^{-(k+1)}.
\]
For the high tail fix $h=(a-b)/2$. Integrating $e^{bv}\anaone_{v>Lt}\le e^{-hLt}e^{(b+h)v}$
gives
\[
 \int_{v>Lt}\abs{E_j}\,\dd\mu_k\le Be^{-hLt}\Bigl(\frac{a}{a-b-h}\Bigr)^{k+1}.
\]
For $k=\ceil{\kappa t}$ choose $\kappa$ with $\kappa\log(R/a)>\gamma+2$ and
$\kappa\log2>a+b+\gamma+2$, and then $L>1$ with $hL>\gamma+2+\kappa\log(a/(a-b-h))$. Each of the
three contributions is $O(e^{-(\gamma+2)t})$. The same tail estimates applied to the majorant
$e^{bv}\ge1$ show $\mu_k([t,Lt])\to1$. Normalizing the restricted measure gives the stated
bound. For continuous $E_j$ the vector of means lies in the convex hull of the compact image of
$[t,Lt]$ in $\R^q$, and Carath\'eodory's theorem gives the finite combination.
\end{proof}

\begin{corollary}[Moving coefficients]\label{ana:cor:window}
Under the hypotheses of \cref{ana:thm:window}, for any $A,\eta>0$ there is a fixed $L>1$ such
that for all large $t$ and every vector of real coefficients with $\abs{c_j}\le e^{At}$, some
$u\in[t,Lt]$ satisfies
\[
 \sum_{j=1}^qc_jE_j(u)\le e^{-\eta t}.
\]
If the $E_j$ are continuous, some $u\in[t,Lt]$ satisfies
$\bigl|\sum_jc_jE_j(u)\bigr|\le e^{-\eta t}$. The measure $\nu_t$ does not depend on the
coefficients. The point $u$ may depend on them.
\end{corollary}

\begin{proof}
Apply \cref{ana:thm:window} with $\gamma>A+\eta+1$. The mean of the combination under $\nu_t$
has absolute value at most $qe^{(A-\gamma)t}$, which is eventually smaller than $e^{-\eta t}$.
If every point of the interval gave a value above $e^{-\eta t}$, so would the mean. A continuous
real function on a compact interval attains its mean.
\end{proof}

\begin{remark}[No common point]\label{ana:rem:nocommon}
Small averages do not give a single point at which all $E_j$ are small. For $E_1(v)=\sin v$ and
$E_2(v)=\cos v$ the hypotheses hold, the Laplace poles being at $\pm i$, but
$E_1(v)^2+E_2(v)^2=1$. The vector statement yields a common weighted average, which is useful
only where the arithmetic inequality is stable under averaging.
\end{remark}

Call a real measurable function $E$ on $[0,\infty)$ \emph{polar} if $\abs{E(v)}\le Ce^v$ and
the function
\[
 \Phi(s)=\int_0^\infty E(v)e^{(1-s)v}\,\dd v\qquad(\Re s>2)
\]
continues meromorphically to $\Re s>0$ with at most simple poles and is holomorphic at every
point of $\Re s\ge1$. This
is the class to which the source applies its window lemma. The remainders $v\mapsto R(e^v)$ of
\cref{ana:lem:mellin} belong to it whenever $A$ is meromorphic on $\Re s>0$ with simple poles
and no pole on $\Re s\ge1$ other than $s=1$; in particular the source's two errors do
(\cref{ana:ex:source}). For a polar function the Laplace transform $F(z)=\Phi(z+1)$ is
holomorphic near $\Re z\ge0$, so the hypotheses of \cref{ana:thm:window} hold with $b=1$ and
any $a>1$.

\begin{theorem}[Sharp universal schedule; \stU]\label{ana:thm:schedule}
Let $H(t)\ge t$ for all large $t$, and fix $a,b\ge0$ with $a+b>0$. The following are
equivalent.
\begin{enumerate}[label=\textup{(\roman*)}]
\item For every pair of polar functions $E,E'$, all sufficiently large $t$ and every
 $0\le X\le e^{bt}$, some $u\in[t,H(t)]$ satisfies $E(u)-XE'(u)\le e^{-at}$.
\item $H(t)/t\to\infty$.
\end{enumerate}
\end{theorem}

\begin{proof}
(ii)$\Rightarrow$(i). For a given pair, \cref{ana:cor:window} supplies a constant $L$ such
that a suitable $u$ exists in $[t,Lt]$, and $Lt\le H(t)$ eventually.

(i)$\Rightarrow$(ii). Suppose $H(t_i)\le L_0t_i$ along a sequence $t_i\to\infty$. Choose
$0<\delta<\min\{1,(a+b)/L_0\}$ and take $E=0$, $E'(v)=-e^{-\delta v}$, $X=e^{bt_i}$. The
transform of $E'$ is $-1/(s-1+\delta)$, with one simple pole at $1-\delta$, so $E'$ is polar.
For every $u\in[t_i,H(t_i)]$,
\[
 E(u)-XE'(u)\ge e^{(b-\delta L_0)t_i}>e^{-at_i}.
\]
So (i) fails, even for continuous functions.
\end{proof}

The counterexample concerns the full analytic class. For a fixed pair the window is linear,
with a constant depending on the pair through the radius $R$ in the proof of
\cref{ana:thm:window}. No uniformity in unbounded conductors is claimed, and arithmetic
subfamilies with larger pole-free regions can have uniform constants.

\begin{corollary}[The second scale of the source; \stU]\label{ana:cor:source}
Fix a modulus $M\ge1$ and a prime $d\equiv7\pmod8$ with $d\nmid M$, and let $R_{d,M}$, $R_{d,1}$
be the remainders of \cref{ana:ex:source}. There is $L=L(d,M)>1$ such that for all
sufficiently large $x$ and every $0\le X\le x^{10^6}$ there is $x_1$ with
\[
 x^{10^6}\le x_1\le x^{10^6L},\qquad
 \bigl|R_{d,M}(x_1)-XR_{d,1}(x_1)\bigr|\le x^{-10^6}.
\]
\end{corollary}

\begin{proof}
The two functions $v\mapsto R_{d,M}(e^v)$ and $v\mapsto R_{d,1}(e^v)$ are continuous and polar.
Apply \cref{ana:cor:window} with $A=\eta=1$ and $t=10^6\log x$.
\end{proof}

In the source the coefficient $X$ is a ratio of two cardinalities bounded by $e^t$ with
$t=10^6\log x$, and the admissible upper endpoint for $x_1$ is $\exp((10^6\log x)^2)$. For
large $x$ the new window lies inside that one, so the main inequality of the source and its
exponent are unchanged. The exponent $1/2$ comes from the number of theta lifts and is not
affected by the window. What changes is the range of auxiliary levels that a new geometric
construction would have to support: a fixed power of $x$ instead of a quasi-polynomial range.

\begin{corollary}[Two-scale certificate; \stU]\label{ana:cor:interface}
Let $R$ be a real function on $[1,\infty)$, let $E_1,\dots,E_q$ satisfy the hypotheses of
\cref{ana:thm:window}, and fix $A_0,A_1>0$. Let $L$ be the constant of \cref{ana:cor:window} for
$A=A_1$ and some $\eta$ with $A_0\eta>1/2$. Suppose that for every $\eps>0$ and every
sufficiently large $x$, with $t=A_0\log x$, there are real coefficients $c_j(x)$ with
$\abs{c_j(x)}\le e^{A_1t}$ such that
\begin{equation}\label{ana:eq:interface}
 R(x)\le C_\eps x^{-1/2+\eps}+\sum_{j=1}^qc_j(x)E_j(u)\qquad\text{for every }u\in[t,Lt].
\end{equation}
Then $R(x)\le(C_\eps+1)x^{-1/2+\eps}$ for all sufficiently large $x$.
\end{corollary}

\begin{proof}
The coefficients do not depend on $u$. Choose $u$ by \cref{ana:cor:window}, so that the sum in
\eqref{ana:eq:interface} is at most $e^{-\eta t}=x^{-A_0\eta}\le x^{-1/2}$.
\end{proof}

The functions $E_j$ may be centred prime-only remainders $v\mapsto R_w(e^v)$ of class
functions, the finite differences $v\mapsto\Delta_\lambda(e^v)$ of \cref{ana:sec:fd}, or the
errors of the source: their transforms are holomorphic near $\Re s\ge1$. Combined with
\cref{ana:thm:primeonly} or \cref{ana:thm:fd}, an inequality of the form
\eqref{ana:eq:interface} for the corresponding statistic $R$ would therefore prove the
criterion. It is a sufficient interface for a geometric comparison with several comparison
scales. No such inequality is available for a function outside the scope of
\cref{src:thm:dgrh}. The form of this certificate with the source window asks for
\eqref{ana:eq:interface} on the larger range $x^{L}\le e^u\le\exp(L^2(\log x)^2)$.

\subsection{Thin sets of twists and of scales}\label{ana:sec:sparse}
{\raggedright\sources{grh_cm_prime_applications; grh_transfer_analysis_2026-11-25; grh_lean_transfer_analysis; grh-torsion-and-prime-sampling; grh_method_transfers_research; grh_method_transfers_sharp_torsion}}

The criteria above ask for an inequality at every large real scale, and the criterion of the
source asks for it at every auxiliary prime of a residue class. Both sets of inputs can be
thinned.

\paragraph{Thin sets of twists for the source remainder.}
Fix an integer $Q>0$ divisible by $8$ and a reduced residue class $a\bmod Q$ with
$a\equiv7\pmod8$. For primes $d\equiv a\pmod Q$ let $R_d=R_{d,Q}$ be the remainder of
\cref{ana:ex:source}. The following input is taken from the source (its lemma
\code{zero_density_quadTwist}, with $\gamma=1/2+9\delta$ and exponent $1-\delta/2$; see
\cref{src:sec:mechanism}): for $1/2<\gamma<1$ and $H\ge0$, the number of primes $d\le X$,
$d\equiv a\pmod Q$, for which some $L(s,\xi\chi_d)$ with $\xi\bmod Q$ has a zero in
$\Re s\ge\gamma$, $\abs{\Im s}\le H$, is
\begin{equation}\label{ana:eq:zerodensity}
 O_{Q,a,\gamma,H}\bigl(X^{1-(\gamma-1/2)/18}\bigr).
\end{equation}

\begin{theorem}[Sparse one-sided zero detection; \stU]\label{ana:thm:sparse}
Fix $1/2<\beta<1$. Let $S$ be a set of primes $d\equiv a\pmod Q$ and
$S(X)=\#\{d\in S:d\le X\}$. Suppose that
\begin{equation}\label{ana:eq:S}
 \limsup_{X\to\infty}\frac{\log(1+S(X))}{\log X}\ge1-\frac{\beta-1/2}{18},
\end{equation}
and that for each $d\in S$ there are a constant $C_d>0$ and an increasing unbounded sequence
$x_j\ge1$, both allowed to depend on $d$, with
\begin{equation}\label{ana:eq:E}
 x_{j+1}-x_j=O_d\bigl(x_j^{(1+\beta)/2}\bigr),\qquad
 R_d(x_j)\le C_dx_j^{\beta-1}\quad(j\text{ sufficiently large}).
\end{equation}
Then for every Dirichlet character $\psi$ modulo $Q$, $L(s,\psi)$ is nonzero in $\Re s>\beta$,
the pole of the principal character excepted.
\end{theorem}

\begin{proof}
\emph{Recovery of the intervening scales.} Write $t=1/x$. On a compact interval
$[u,v]\subset(0,\infty)$ the function $F(t)=\sum_{n\ge1}\lambda_{d,Q}(n)(1/n-t)_+$ is a finite
sum of convex functions, hence lies below its chord, and
$R_d(1/t)=F(t)+(\log t)/\varphi(Q)-c_{d,Q}$. Write $If$ for the affine interpolant of $f$ from
$u$ and $v$. Since $t\mapsto\log t+t^2/(2u^2)$ is convex on $[u,v]$,
\[
 0\le\log t-I(\log)(t)\le\frac{(t-u)(v-t)}{2u^2}\le\frac{(v-u)^2}{8u^2}.
\]
Take $u=1/x_{j+1}$ and $v=1/x_j$. Then $(v-u)/u=(x_{j+1}-x_j)/x_j$, so the interpolation loss
is $O_d(x_j^{\beta-1})$, and $x_{j+1}/x_j\to1$. With \eqref{ana:eq:E} at the two endpoints this
gives $R_d(x)\le K_dx^{\beta-1}$ for all large $x$, for some $K_d>0$.

\emph{Selection of a twist.} Suppose $L(\rho,\psi)=0$ with $\sigma=\Re\rho>\beta$. The
classical nonvanishing on $\Re s\ge1$ leaves $\sigma<1$. Choose $\gamma$ with
$\beta<\gamma<\sigma$ and put $H=\abs{\Im\rho}$. Call $d$ bad if some $L(s,\xi\chi_d)$,
$\xi\bmod Q$, has a zero in $\Re s\ge\gamma$, $\abs{\Im s}\le H$. By
\eqref{ana:eq:zerodensity} the bad primes up to $X$ number $O(X^{1-(\gamma-1/2)/18})$, and this
exponent is strictly smaller than the right side of \eqref{ana:eq:S}. Hence some $d\in S$ is
not bad. Fix it. Since $K_dx^{\beta-\gamma}\to0$, we have $R_d(x)\le x^{\gamma-1}$ eventually.

\emph{The fixed twist.} In \eqref{ana:eq:srcres} the twisted orders vanish in the box. On the
real segment $\gamma<s<1$ the residues of $\Phi_{d,Q}$ are therefore nonnegative, and for real
$s\ge1$ there is no pole. \Cref{ana:thm:onesided} with \cref{ana:cor:forms}(a) shows that
$\Phi_{d,Q}$ has no pole in $\Re s>\gamma$. But its residue at $\rho$ is the total multiplicity
of $\rho$ as a zero of the untwisted functions, which is positive, divided by
$\varphi(Q)\rho(1-\rho)\ne0$. This is a contradiction.
\end{proof}

The order of quantifiers matters: first $\rho$, then a successful $d$, then the eventual
estimate for that $d$. No uniformity in the constants $C_d$ or in the thresholds is required.

\begin{corollary}\label{ana:cor:density}
Condition \eqref{ana:eq:S} holds in each of the following cases:
\begin{enumerate}[label=\textup{(\alph*)}]
\item $S$ has positive upper relative density among the primes $d\equiv a\pmod Q$;
\item $\limsup_{X\to\infty}\log(1+S(X))/\log X=1$;
\item $\sum_{d\in S}(\log d)/d=+\infty$;
\item $S(X)/X^{1-(\beta-1/2)/18}$ is unbounded.
\end{enumerate}
In particular, if for every $d$ in such a set one has $R_d(x)\le C_dx^{\beta-1}$ for all
sufficiently large real $x$, then no $L(s,\psi)$, $\psi\bmod Q$, vanishes in $\Re s>\beta$.
A single set $S$ with counting exponent $1$, on which the bound holds for every $\beta>1/2$,
excludes all zeros to the right of $1/2$.
\end{corollary}

\begin{proof}
(a) follows from the prime number theorem in the progression and implies (b). (b) and (d) are
immediate. For (c), if \eqref{ana:eq:S} failed then $S(X)\le X^{1-\eta}$ for some $\eta>0$ and
all large $X$, and the sum would converge by partial summation.
\end{proof}

\begin{remark}\label{ana:rem:schedule}
(1) The grid $x_j=j^4$ satisfies the gap condition in \eqref{ana:eq:E} for every
$\beta\ge1/2$. The fourth-power mesh and the chord interpolation for the prime sum of the
ordinary zeta function are due to Setiawan \cite{Setiawan2026Primacohedron}. The exact
interpolation loss and its sharpness for general positive measures are given in
\cref{sec:nab}.

(2) If a zero $\rho$ with $\Re\rho>\beta$ existed, then for every $\gamma\in(\beta,\Re\rho)$
the number of admissible primes $d\le X$ with
$\limsup_{j\to\infty}j^{4(1-\beta)}R_d(j^4)<+\infty$ would be
$O(X^{1-(\gamma-1/2)/18})$, since every such $d$ is bad for one fixed box. So the normalized
discrepancy at fourth powers would have arbitrarily large positive excursions for all
admissible primes outside a set of power-saving size.

(3) An explicit thin schedule is the following. Let $p_1<p_2<\cdots$ be the primes
$\equiv a\pmod Q$ and $d_k=p_{\floor{k\log^2(k+2)}}$. Then $d_k\sim\varphi(Q)k\log^3k$, the
reciprocals of the $d_k$ have a convergent sum, and the counting exponent of $\{d_k\}$ is $1$.
If for every $k$ and every $\eps>0$ one has
$R_{d_k}(j^4)\le C_{k,\eps}j^{-2+4\eps}$ for all large $j$, then no $L(s,\psi)$ modulo $Q$
vanishes to the right of $1/2$. For an arbitrary modulus $q$ take $Q=18480q$ and $a=Q/2-1$;
the nonvanishing descends to $q$, because changing the level changes finitely many Euler
factors. For these parameters the estimate proved in the source implies the sampled
inequalities. The statement describes the weaker input that an alternative proof would have
to supply. Bounds for finitely many twists, or at finitely many scales, give nothing.
\end{remark}

\paragraph{Prescribed sparse scales for a Dedekind zeta function.}
Let $K$ be a number field and $R_K$ the remainder \eqref{ana:eq:RK}. All sums over the
nontrivial zeros $\rho$ of $\zeta_K$ count multiplicity. Put
$\Theta_K=\sup_\rho\Re\rho\in[1/2,1]$ and $c_\rho=1/\bigl(\rho(1-\rho)\bigr)$. The smoothed
explicit formula is
\begin{equation}\label{ana:eq:explicit}
 R_K(x)=\sum_\rho c_\rho\,x^{\rho-1}+O_K\bigl(x^{-1}\log(2x)\bigr)\qquad(x\ge2).
\end{equation}
It is obtained by shifting the inverse Mellin integral of the uncentred ramp from $\Re s=2$
to $\Re s=-1/2$ along heights avoiding zeros. The terms at $s=1$ give the centring, the pole
of order at most two at $s=0$ gives $O_K(x^{-1}\log(2x))$, and the new line contributes
$O_K(x^{-3/2})$ by the functional equation and Stirling's formula; compare
\cite[\S\S1.2--1.3]{IharaEulerKronecker}. The zero count $N_K(T)=O_K(T\log(2T))$ gives
$\sum_\rho\abs{c_\rho}<\infty$, and by partial summation
\begin{equation}\label{ana:eq:zerosums}
 \sum_{\abs{\Im\rho}\le T}\frac1{\abs\rho}\ll_K\log^2(2T),\qquad
 \sum_{\abs{\Im\rho}>T}\frac1{\abs{\rho(1-\rho)}}\ll_K\frac{\log(2T)}T\qquad(T\ge1).
\end{equation}

\begin{theorem}[Prescribed grids; \stU]\label{ana:thm:powers}
Let $y_j$ be an unbounded increasing sequence with $\log(y_{j+1}/y_j)=O(y_j^{-\delta})$ for
some $\delta>0$; for instance $y_m=m^k$ for a fixed integer $k\ge1$, with $\delta=1/k$.
\begin{enumerate}[label=\textup{(\roman*)}]
\item For every $0<\beta<\Theta_K$,
 $\limsup_{j\to\infty}y_j^{1-\beta}R_K(y_j)=+\infty$.
\item For every $\beta>\Theta_K$, $x^{1-\beta}R_K(x)\to0$ as $x\to\infty$.
\item If all nontrivial zeros of $\zeta_K$ lie on the critical line, then
 $\abs{R_K(x)}\ll_Kx^{-1/2}$.
\end{enumerate}
In particular, for every fixed $k\ge1$ the Riemann hypothesis for $\zeta_K$ is equivalent to
an eventual one-sided bound $R_K(m^k)\le C_{K,k}\,m^{-k/2}$ at the integers $m$.
\end{theorem}

\begin{proof}
(ii) and (iii) follow from \eqref{ana:eq:explicit} and the absolute convergence of
$\sum_\rho c_\rho$. For (i), let $0<h\le1$. Since
$\abs{e^{(\rho-1)h}-1}\le\min(2,\abs{\rho-1}h)$, splitting the sum in \eqref{ana:eq:explicit}
at $\abs{\Im\rho}=1/h$ and using \eqref{ana:eq:zerosums} gives
\[
 \abs{R_K(xe^h)-R_K(x)}\ll_Kx^{\Theta_K-1}h\log^2(2/h)+x^{-1}\log(2x).
\]
Suppose that the upper limit in (i) is finite for some $\beta<\Theta_K$, so that
$R_K(y_j)\le Cy_j^{\beta-1}$ for all $j$. Choose $\tau$ with
$\max(\beta,\Theta_K-\delta)<\tau<\Theta_K$. For $y_j\le x\le y_{j+1}$ the variation estimate
gives $R_K(x)\le C'x^{\tau-1}$ on the whole ray. By \cref{ana:thm:zerotest} and
\cref{ana:ex:instances}(a), $\zeta_K$ has no zero in $\Re s>\tau$, which contradicts the
definition of $\Theta_K$. The last assertion follows from (i) at $\beta=1/2$ and from (iii).
\end{proof}

Part (i) holds even if no zero attains the supremum $\Theta_K$. The grid may be far sparser
than the one allowed by the convexity argument in the proof of \cref{ana:thm:sparse}, which
at $\beta=1/2$ reaches fourth powers; the price is the use of the explicit formula, so the
statement is special to a function all of whose residues have one sign.

\begin{corollary}[A finite statistic on power grids; \stU]\label{ana:cor:finite}
Fix $k\ge1$ and put $D_k(m)=R_K(m^k)-R_K((2m)^k)$, which is the difference of the two finite
ramp sums plus $k\log2$. The Riemann hypothesis for $\zeta_K$ is equivalent to
$\sup_{m\ge1}m^{k/2}D_k(m)<\infty$.
\end{corollary}

\begin{proof}
Under the Riemann hypothesis use \cref{ana:thm:powers}(iii). Conversely $R_K(x)\to0$
unconditionally, by \eqref{ana:eq:explicit}, dominated convergence and the nonvanishing on
$\Re s=1$. Telescoping over $m,2m,4m,\dots$ gives $R_K(m^k)\le Cm^{-k/2}/(1-2^{-k/2})$ from a
nonnegative upper bound $C$ for the statistic, and \cref{ana:thm:powers} applies.
\end{proof}

\begin{proposition}[One base for every number field; \stU]\label{ana:prop:generic}
Let $0<\lambda_n\uparrow\infty$ with $\lambda_{n+1}/\lambda_n\to1$. There is a dense countable
intersection of open sets of bases $b>1$ for which conclusions \textup{(i)} and
\textup{(ii)} of \cref{ana:thm:powers} hold on the grid $x=b^{\lambda_n}$, for every number
field $K$ simultaneously. For instance $\lambda_n=\exp(n/\log(n+2))$ uses asymptotically
$(\log\log X)\log\log\log X$ samples up to $X$ for fixed $b$.
\end{proposition}

\begin{proof}
For a continuous function $F$ and a rational $h<\limsup_{x\to\infty}F(x)$, the set of $b$ with
$F(b^{\lambda_n})>h$ for some $n\ge N$ is open, and it is dense because for $1<a<c$ the
intervals $(a^{\lambda_n},c^{\lambda_n})$ eventually overlap and cover a ray. Intersect over
$h$, over $N$, and over the countably many functions $F(x)=x^{1-\beta}R_K(x)$ with rational
$0<\beta<1$. By \cref{ana:thm:powers} with $y_j=j$ the upper limit of $F$ on the whole ray is
$+\infty$ for $\beta<\Theta_K$, and intermediate rational exponents give the real ones. This
is a Croft--Kingman sampling argument \cite{BinghamOstaszewskiKingman}.
\end{proof}

The generic statement leaves open the status of a prescribed base. The next theorem treats
every base on one explicit set of integer exponents. Put
\[
 M_j=2^{2^{j^2}},\qquad\anaE=\bigcup_{j\ge1}\{M_j,2M_j,\dots,jM_j\}.
\]
For every $b>1$ at most $O_b(1+\log\log\log X)$ of the cutoffs $b^e$, $e\in\anaE$, lie below
$X$: indeed $b^{M_j}\le X$ implies $j^2=O_b(1+\log\log\log X)$, and the first $j$ blocks
contain $j(j+1)/2$ terms.

\begin{theorem}[Extremal values on a prescribed sparse set; \stU]\label{ana:thm:blocks}
Let $K$ be a number field with signature $(r_1,r_2)$ and discriminant $d_K$, and let $b>1$.
\begin{enumerate}[label=\textup{(\roman*)}]
\item If every nontrivial zero of $\zeta_K$ has real part $1/2$, then
 \[
  \lim_{j\to\infty}\max_{1\le k\le j}b^{kM_j/2}R_K\bigl(b^{kM_j}\bigr)
  =\sum_\rho\frac1{1/4+(\Im\rho)^2}=:S_K ,
 \]
 and
 $S_K=2+2\gamma_K+\log\abs{d_K}-r_1(\gamma+\log4\pi)-2r_2(\gamma+\log2\pi)$, where $\gamma$ is
 Euler's constant. Moreover, for every $\eps>0$ a positive proportion of the indices $k$ in
 every sufficiently large block gives a value within $\eps$ of $S_K$.
\item Unconditionally, the upper limit of $x^{1-\Theta_K}R_K(x)$ over $x=b^e$, $e\in\anaE$, is
 strictly positive if and only if some zero has real part $\Theta_K$. If no zero attains
 $\Theta_K$, then $x^{1-\Theta_K}R_K(x)\to0$ on the whole ray.
\end{enumerate}
\end{theorem}

By \cref{src:thm:dgrh} the hypothesis of (i) holds for $K=\Q$ and for every abelian field, and
$S_\Q=2+\gamma-\log(4\pi)=0.0461914\ldots$.

\begin{proof}
Write $\Theta=\Theta_K$ and $\rho=\sigma+i\gamma_\rho$. Multiply \eqref{ana:eq:explicit} by
$x^{1-\Theta}$. By dominated convergence the sum over $\Re\rho<\Theta$ tends to zero, and so
does the error term. The only surviving sum is
$\sum_{\Re\rho=\Theta}c_\rho\exp(i\gamma_\rho\log x)$, which is empty if $\Theta$ is not
attained.

Fix a finite list of $n$ of its ordinates, set $\theta_i=\gamma_i\log b/(2\pi)$, and take an
integer $q\ge2$. The fractional parts of $0,M_j\theta,\dots,q^nM_j\theta$ occupy $q^n$
half-open boxes of side $1/q$ in $[0,1)^n$. Two of these points share a box, and their
difference supplies $1\le k\le q^n$ with $\norm{kM_j\theta_i}_{\R/\Z}\le1/q$ for all $i$. For
$j\ge q^n$ this $k$ belongs to the $j$th block. In fact a positive proportion of each large
block works: for $N=j+1$ points and $B=q^n$ boxes, the number of same-box pairs is at least
$(N^2/B-N)/2$ by the Cauchy--Schwarz inequality, each positive difference occurs in at most
$N$ pairs, and so there are at least $(N/B-1)/2$ suitable values of $k$. Absolute summability
allows one to discard a tail of arbitrarily small norm and then to choose $q$ so that all
remaining phases are arbitrarily close to $1$. Hence a subsequence of the normalized values
tends to $\sum_{\Re\rho=\Theta}c_\rho$. This sum is real by conjugation, and each summand has
positive real part because $\Re[\rho(1-\rho)]=\sigma(1-\sigma)+\gamma_\rho^2>0$. This proves
(ii).

On the critical line every $c_\rho=1/(1/4+\gamma_\rho^2)$ is positive, and the triangle
inequality gives the matching upper bound on every block. The value of $S_K$ in terms of the
invariants of $K$ follows from \cite[(1.2.2), (1.2.3), (1.3.11)]{IharaEulerKronecker}.
\end{proof}

For an arbitrary field, a finite upper bound for $\sqrt x\,R_K(x)$ on this grid therefore
excludes an attained rightmost zero with real part greater than $1/2$: on the subsequence
constructed in the proof the factor $x^{\Theta-1/2}$ tends to infinity. The sparsity counts
cutoffs. It gives no saving in the evaluation of a prime sum at such a cutoff.

\begin{example}\label{ana:ex:cuberoot}
For $K=\Q(\sqrt[3]{2})$, of discriminant $-108$, the weights are computable: for $p\nmid6$ the
degrees $f_i$ of the irreducible factors of $T^3-2$ modulo $p$ give
$\Lambda_K(p^j)=(\log p)\sum_{f_i\mid j}f_i$, and $\Lambda_K(p^j)=\log p$ for $p\in\{2,3\}$,
where $p$ is totally ramified. Here $\zeta_K=\zeta\cdot L(s,\mathrm{std})$ with the
two-dimensional standard representation of $S_3$, which is outside the scope of
\cref{src:thm:dgrh}. The sampled bounds of \cref{ana:thm:powers,ana:cor:finite} for this field
are open.
\end{example}

\paragraph{Simple zeros.}
The same remainder detects multiplicity on a sparse sequence. Let
\[
 x_n=2^n3^{n^2}.
\]
Only $O(\sqrt{\log X})$ of these cutoffs lie below $X$. Assume that all nontrivial zeros of
$\zeta_K$ lie on the critical line. For real $t$ let $m_K(t)$ be the multiplicity of $1/2+it$
as a zero, and let $\Gamma_K$ be the set of distinct ordinates. By \eqref{ana:eq:explicit},
\begin{equation}\label{ana:eq:apseries}
 \sqrt x\,R_K(x)=\sum_{t\in\Gamma_K}a_tx^{it}+o(1),\qquad
 a_t=\frac{m_K(t)}{1/4+t^2},\qquad\sum_ta_t<\infty,
\end{equation}
with uniform convergence in $\log x$. Let $\xi_K(s)=s(s-1)\abs{d_K}^{s/2}
\Gamma_\R(s)^{r_1}\Gamma_\C(s)^{r_2}\zeta_K(s)$ be the completed entire zeta function, with
$\Gamma_\R(s)=\pi^{-s/2}\Gamma(s/2)$ and $\Gamma_\C(s)=2(2\pi)^{-s}\Gamma(s)$, and
$H_K=(\log\xi_K)'$.

\begin{theorem}[Multiplicities from a sparse sequence; \stG\ when $K$ is abelian]\label{ana:thm:simple}
Let $K$ be a number field all of whose nontrivial zeros lie on the critical line.
\begin{enumerate}[label=\textup{(\roman*)}]
\item For every real $u$,
 \[
  \lim_{N\to\infty}\frac1N\sum_{n=1}^N\sqrt{x_n}\,R_K(x_n)\,x_n^{-iu}=\frac{m_K(u)}{1/4+u^2}.
 \]
\item The limit
 \[
  \lim_{N\to\infty}\frac1N\sum_{n=1}^Nx_nR_K(x_n)^2
  =\sum_{t\in\Gamma_K}\frac{m_K(t)^2}{(1/4+t^2)^2}
 \]
 exists, and it exceeds $4H_K(1)-2H_K'(1)$ by exactly
 $\sum_{t}m_K(t)\bigl(m_K(t)-1\bigr)/(1/4+t^2)^2\ge0$.
\item \textup{(\stG.)} For $K=\Q$, all nontrivial zeros of $\zeta$ are simple if and only if
 \[
  \lim_{N\to\infty}\frac1N\sum_{n=1}^Nx_nR_\Q(x_n)^2
  =6+2\gamma-2\log(4\pi)+4\gamma_1+2\gamma^2-\frac{\pi^2}4\approx0.00007420127287493 ,
 \]
 where $\zeta(1+u)=u^{-1}+\gamma-\gamma_1u+O(u^2)$. The excess $D$ of the limit over this
 constant satisfies $0\le D<1.67\cdot10^{-34}$.
\end{enumerate}
\end{theorem}

The hypothesis holds for $K=\Q$ and for abelian fields by \cref{src:thm:dgrh}; for a general
field it is an additional assumption. No assumption is made about rational relations among
the ordinates.

\begin{proof}
(i) For every real $u$,
\[
 \lim_{N\to\infty}\frac1N\sum_{n=1}^Nx_n^{iu}=\begin{cases}1,&u=0,\\0,&u\ne0.\end{cases}
\]
Indeed, for $u\ne0$ at least one of $u\log2/(2\pi)$ and $u\log3/(2\pi)$ is irrational, since
$\log2/\log3$ is irrational, and Weyl's theorem applies to the polynomial with these linear
and quadratic coefficients \cite{Weyl1916Gleichverteilung}. Multiply \eqref{ana:eq:apseries}
by $x_n^{-iu}$ and average termwise: absolute convergence controls the tail uniformly in $N$,
and the average of the $o(1)$ term tends to zero.

(ii) The prime sum is real, so its square is its squared absolute value. The double series
obtained from \eqref{ana:eq:apseries} converges absolutely, and by the same averaging only
pairs with equal ordinates survive. The Hadamard product of genus one and the symmetry
$\xi_K(s)=\xi_K(1-s)$ give
\[
 \sum_\rho\frac1{\rho(1-\rho)}=H_K(1)-H_K(0)=2H_K(1),\qquad
 \sum_\rho\frac1{(1-\rho)^2}=-H_K'(1).
\]
With $[z(1-z)]^{-2}=z^{-2}+(1-z)^{-2}+2/[z(1-z)]$ and the same symmetry,
\[
 \sum_{t\in\Gamma_K}\frac{m_K(t)}{(1/4+t^2)^2}=4H_K(1)-2H_K'(1),
\]
and subtracting this from the limit gives the stated excess.

(iii) For $K=\Q$ the expansion at $s=1$ gives
\[
 H_\Q(1)=1+\tfrac12\gamma-\tfrac12\log(4\pi),\qquad
 H_\Q'(1)=-1-2\gamma_1-\gamma^2+\frac{\pi^2}8,
\]
which is the stated constant, and $D$ vanishes exactly when all zeros are simple. For the
enclosure let $T_0=3\cdot10^{12}$. The zeros up to height $T_0$ are simple
\cite{PlattTrudgian2021Riemann}. For $t\ge T_0$ the zero-counting function satisfies
$\abs{N(t)-(t/2\pi)\log(t/(2\pi e))}\le0.1038\log t+0.2573\log\log t+9.3675<\log t$
\cite[Corollary~1.2]{HasanalizadeShenWong2021Counting}. Hence $N(t)\le t\log t$ there and a
jump of $N$ at $t>T_0$ has size at most $2\log t$. By conjugation symmetry and Stieltjes
integration by parts, in which the boundary term discarded at $T_0$ is nonpositive,
\[
 D\le4\int_{(T_0,\infty)}\frac{\log t}{t^4}\,\dd N(t)
 \le16\int_{T_0}^\infty\frac{\log^2t}{t^4}\,\dd t
 =\frac{16}{T_0^3}\Bigl(\frac{\log^2T_0}3+\frac{2\log T_0}9+\frac2{27}\Bigr)<1.67\cdot10^{-34}.
\]
The rational evaluation of the last bound is recorded in the file
\begin{center}
\artifact{grh-torsion-and-prime-sampling/verification.txt}.
\end{center}
\end{proof}

The identity in (ii) has the same shape for the polynomial $n^d$ in place of $n^2$, for each
fixed $d\ge2$. Continuous mean squares of prime errors also detect repeated zeros
\cite{Caceres2026Constant}. Here the reciprocal weights make the series
\eqref{ana:eq:apseries} absolutely convergent, and the polynomial choice of $\log x_n$ allows
sampling on one sparse sequence. The arithmetic question that remains is an upper bound for
the sampled mean square by the explicit constant in (iii). It would require control of the
off-diagonal terms when the prime side of the mean square is expanded at these cutoffs.

\paragraph{Negative excursions of the source remainder.}
The remainder of the source is bounded from above by the source itself. Its excursions below
zero at modulus one can be located as follows. Fix a prime $d\equiv7\pmod8$, let
$\chi=\chi_d$, and write $R_d=R_{d,1}$, so that
\[
 R_d(x)=\sum_{n\le x}\Lambda(n)\bigl(1-\chi(n)\bigr)\Bigl(\frac1n-\frac1x\Bigr)
 -\log x+1+\gamma-\frac{L'(1,\chi)}{L(1,\chi)} .
\]
Put
\[
 B_0=2+\gamma-\log(4\pi)=0.0461914179\ldots,\qquad
 B_d=\log\frac d\pi-\gamma+2\,\frac{L'(1,\chi)}{L(1,\chi)} .
\]

\begin{theorem}[Excursion bounds; \stG]\label{ana:thm:excursion}
For every prime $d\equiv7\pmod8$,
\[
 -B_d-B_0\le\liminf_{x\to\infty}\sqrt x\,R_d(x)\le-B_d+B_0 .
\]
A single sequence $x_j\to\infty$ realizes $\sqrt{x_j}\,R_d(x_j)\to B_0-B_d$ for every such $d$
simultaneously, and there is a single sequence of integer exponents $n_j\to\infty$ with
$\sqrt{a^{n_j}}\,R_d(a^{n_j})\to B_0-B_d$ for every integer $a\ge2$ and every $d$. In
particular the lower limit divided by $\log d$ tends to $-1$ as $d\to\infty$.
\end{theorem}

The limits are taken in this order, first in $x$ for fixed $d$. Nothing is asserted about the
first $x$ at which a given excursion occurs.

\begin{proof}
By \cref{src:thm:dgrh} the nontrivial zeros of $\zeta$ and of $L(s,\chi)$ lie on the critical
line. For $F=\zeta$ or $F=L(\cdot,\chi)$ put
\[
 H_F(u)=\sum_{\rho=1/2+i\gamma_\rho}\frac{e^{i\gamma_\rho u}}{1/4+\gamma_\rho^2},
\]
over all nontrivial zeros with multiplicity, central zeros included. The series converges
absolutely and uniformly, and is real. For fixed $d$,
\begin{equation}\label{ana:eq:O1}
 \sqrt x\,R_d(x)=H_\zeta(\log x)-H_{L(\cdot,\chi)}(\log x)+O_d\bigl(x^{-1/2}\bigr).
\end{equation}
To see this, integrate $x^sA(s)/(s(s-1))$ over $\Re s=2$, with
$A=-\zeta'/\zeta+L'/L(\cdot,\chi)$. Mellin inversion gives
$\sum_{n\le x}\Lambda(n)(1-\chi(n))(x/n-1)$. At $s=1$ the residue is
$x\bigl(\log x-1-\gamma+L'(1,\chi)/L(1,\chi)\bigr)$, and at a zero $\rho$ it is
$x^\rho\bigl(m_\zeta(\rho)-m_L(\rho)\bigr)/\bigl(\rho(1-\rho)\bigr)$, coincident zeros
included. Shift the contour to $\Re s=-1/2$. The kernel has a simple pole at $0$, since
$\zeta(0)L(0,\chi)\ne0$ for this odd primitive character, with a residue depending only on
$d$. On the new line the logarithmic derivatives are $O_d(\log(\abs{\Im s}+2))$ by the
functional equations, and the horizontal edges can be sent to infinity along heights at
distance $\gg_d1/\log T$ from both zero sets. The new integral is $O_d(x^{-1/2})$. Dividing
by $\sqrt x$ gives \eqref{ana:eq:O1}.

For the completed entire function $\Xi_F$ the Hadamard product gives
\[
 \frac{\Xi_F'}{\Xi_F}(1)-\frac{\Xi_F'}{\Xi_F}(0)=\sum_\rho\frac1{\rho(1-\rho)}=H_F(0),
\]
and by the functional equation the left side is twice the logarithmic derivative at $1$.
Inserting the gamma factors gives $H_\zeta(0)=B_0$ and $H_{L(\cdot,\chi)}(0)=B_d$. As
$\abs{H_F(u)}\le H_F(0)$, \eqref{ana:eq:O1} proves the lower bound.

For finitely many real ordinates, simultaneous approximation returns all their phases
arbitrarily close to $1$ at arbitrarily large times, whatever their rational relations: the
integer iterates in the finite torus have a convergent subsequence, and differences of
increasingly separated indices return to the identity. Enumerate the ordinates of $\zeta$ and
of the countably many characters, and take $u_j>j$ approximating the first $j$ phases to
within $1/j$. Absolute convergence of each series gives
$H_\zeta(u_j)-H_{L(\cdot,\chi)}(u_j)\to B_0-B_d$ for every $d$, which proves the upper bound
and the simultaneous assertion with $x_j=e^{u_j}$. The same argument with the countably many
frequencies $\gamma_\rho\log a$, $a\ge2$, gives the integer exponents $n_j$. Finally
$L'(1,\chi)/L(1,\chi)=O(\log\log d)$ under \cref{src:thm:dgrh}
\cite[Corollary~3.3.2]{IharaMurtyShimura2009Logarithmic}, which gives $B_d\sim\log d$.
\end{proof}

\begin{proposition}[When the lower endpoint is attained; \stG]\label{ana:prop:parity}
The value $-B_d-B_0$ is a limiting value of $\sqrt x\,R_d(x)$ if and only if every finite
integer relation among the positive ordinates of $\zeta$ and of $L(\cdot,\chi)$ has an even
total coefficient on the ordinates of $\zeta$. Here each occurrence of an ordinate is
labelled by its function, so a positive ordinate shared by the two functions already violates
the condition.
\end{proposition}

\begin{proof}
Equality forces each cosine with positive weight in $H_\zeta$ to approach $-1$ and each
cosine in $H_{L(\cdot,\chi)}$ to approach $+1$, because their individually nonnegative
deficits sum to zero in the limit. The function $\zeta$ has no central zero, and central
zeros of $L$ contribute with the required sign. On each finite set of frequencies,
Kronecker's theorem identifies the closure of the real flow with the subtorus cut out by the
integer relations, and the required sign vector lies in that closure exactly under the
stated parity condition. A diagonal sequence and the summable tails give the sufficiency for
all ordinates.
\end{proof}

For $d=7$ the two endpoints are approximately $-0.30142740$ and $-0.20904456$, and for
$d=23$ they are $-0.61549819$ and $-0.52311535$. Further values are in the file
\begin{center}
\artifact{grh_method_transfers_sharp_torsion/checks/oscillation_bounds.txt}.
\end{center}
These numbers evaluate the proven bounds through
\[
 \frac{L'(1,\chi)}{L(1,\chi)}=\log(2\pi)+\gamma-\frac1{L(0,\chi)}\sum_{a=1}^{d-1}\chi(a)\log\Gamma(a/d),
 \qquad L(0,\chi)=-\frac1d\sum_{a=1}^{d-1}a\chi(a).
\]
They are not estimates of a lower limit by sampling.

\subsection{Sparse Li coefficients}\label{ana:sec:li}
{\raggedright\sources{grh_method_transfers_research}}

\Cref{ana:thm:blocks}(ii) does not see zeros off the line whose real parts approach an
unattained supremum. The block argument gives a full equivalence for a different sequence.
For a number field $K$ put $w_\rho=1-1/\rho$ and define the Li coefficients
\[
 \lambda_n(K)=\lim_{T\to\infty}\sum_{\abs{\Im\rho}\le T}\bigl(1-w_\rho^n\bigr)\qquad(n\ge1).
\]
They are real, and
$\lambda_n(K)=\frac1{(n-1)!}\bigl(\frac{\dd}{\dd s}\bigr)^n\bigl[s^{n-1}\log\xi_K(s)\bigr]_{s=1}$.
Li's criterion for number fields and the dominant exponential terms when it fails are
classical \cite{BombieriLagariasComplements,LagariasLi}. The additional content below is the
prescribed sparse set of indices $\anaE$ of \cref{ana:thm:blocks}.

\begin{theorem}[Sparse Li criterion; \stU]\label{ana:thm:li}
For every number field $K$ the following are equivalent.
\begin{enumerate}[label=\textup{(\roman*)}]
\item All nontrivial zeros of $\zeta_K$ lie on the critical line.
\item $\lambda_n(K)\ge0$ for all sufficiently large $n\in\anaE$.
\item The proportion of indices $k\in\{1,\dots,j\}$ with $\lambda_{kM_j}(K)<0$ tends to zero as
 $j\to\infty$.
\end{enumerate}
If \textup{(i)} fails, let $\varrho>1$ be the largest value of $\abs{w_\rho}$ and $m$ the number
of zeros attaining it, with multiplicity. Then
\[
 \lim_{j\to\infty}\min_{1\le k\le j}\varrho^{-kM_j}\lambda_{kM_j}(K)=-m .
\]
The set $\anaE$ contains $O(1+\log\log N)$ indices up to $N$.
\end{theorem}

\begin{proof}
On the critical line $\abs{w_\rho}=1$ and every $1-\Re(w_\rho^n)$ is nonnegative, so (i) implies
(ii) and (iii). Suppose (i) fails. By the symmetry about the critical line there is a zero with
real part below $1/2$, and then $\abs{w_\rho}>1$, since
$\abs{w_\rho}^2=1+(1-2\Re\rho)/\abs\rho^2$. As $w_\rho\to1$ with increasing height, the largest
modulus $\varrho$ is attained at finitely many zeros, and all other moduli are at most some $r$
with $1<r<\varrho$. For these other zeros Taylor's formula on the segment from $1$ to $w_\rho$
gives, for $n\ge2$,
\[
 \Bigl|1-w_\rho^n-\frac n\rho\Bigr|\le\frac{n(n-1)}{2\abs\rho^2}\,r^{n-2}.
\]
The sum of $\abs\rho^{-2}$ is finite, and the symmetric sum of $1/\rho$ exists. Hence
\[
 \varrho^{-n}\lambda_n(K)=-\sum_{\abs{w_\rho}=\varrho}\Bigl(\frac{w_\rho}{\varrho}\Bigr)^n+o(1).
\]
Apply the same-box argument of the proof of \cref{ana:thm:blocks} to the arguments of the $m$
outermost points. In every sufficiently large block the real parts of their normalized powers
can all be made arbitrarily close to $1$, for a positive proportion of the indices $k$, and the
real part of their sum never exceeds $m$. This gives the limit $-m$, and it contradicts (ii)
and (iii). Counting the blocks with $M_j\le N$ gives the sparsity.
\end{proof}

The proof allows any sequence of positive integer dilations $M_j\to\infty$ with block lengths
tending to infinity. For a field such as $\Q(\sqrt[3]{2})$ the sampled inequalities in (ii)
are open.

For two number fields $K$ and $L$ let $d_n=\lambda_n(K)-\lambda_n(L)$. The next result needs
no hypothesis on the location of zeros.

\begin{theorem}[Arithmetic equivalence from sparse Li coefficients; \stU]\label{ana:thm:rigidity}
If $\zeta_K\ne\zeta_L$, then
\[
 \lim_{H\to\infty}\ \inf_{A\ge1}\ \frac1H\sum_{n=A}^{A+H-1}\abs{d_n}=+\infty\qquad(A,H\in\N).
\]
Thus a bounded mean absolute discrepancy on a single sequence of windows of increasing length
forces $\zeta_K=\zeta_L$, and then $d_n=0$ for every $n$. The windows may be arbitrarily far
apart: those starting at $2^{2^{j^2}}$ and of length $j$ use only $O(1+\log\log N)$ indices up
to $N$. The same holds with $\abs{d_n}^2$ in place of $\abs{d_n}$.
\end{theorem}

\begin{proof}
The statement for squares follows from the one for absolute values by the Cauchy--Schwarz
inequality. Write $m_w$ for the difference of the multiplicities, for $K$ and for $L$, of the
transformed zero $w=1-1/\rho$. Second differences of the two symmetric sums give
\[
 a_n:=-(d_{n+2}-2d_{n+1}+d_n)=\sum_wc_ww^n,\qquad c_w=m_w(1-w)^2,\qquad\sum_w\abs{c_w}<\infty,
\]
because $\abs{1-w}=\abs\rho^{-1}$. All $\abs w$ are bounded. Suppose that there are windows
$[A,A+H-1]$ with $H\to\infty$ on which the mean of $\abs{d_n}$ is at most $C$. Then
$\abs{d_n}\le CH$ on each window.

\emph{Points off the unit circle.} Suppose some $w$ with $m_w\ne0$ has modulus above $1$, and
let $\varrho$ be the largest such modulus. It is attained, at finitely many distinct points
$w_1,\dots,w_t$, and the other moduli with $m_w\ne0$ are at most some $r<\varrho$. Put
$z_i=w_i/\varrho$. For $N=A+H-t-2$ and $0\le h<t$,
\[
 \varrho^{-N-h}a_{N+h}=\sum_{i=1}^tz_i^h\,c_{w_i}z_i^N+O\bigl((r/\varrho)^{N+h}\bigr).
\]
The left sides tend to zero, since $N\ge H-t-1$ and $\abs{a_{N+h}}\le4CH$. The Vandermonde
matrix $(z_i^h)_{h,i}$ is invertible, so $c_{w_i}z_i^N\to0$, which contradicts
$\abs{c_{w_i}}>0$. Hence $m_w=0$ for $\abs w>1$. The reflection $\rho\mapsto1-\rho$ induces
$w\mapsto1/w$ and preserves multiplicity, so $m_w=0$ for $\abs w<1$ as well.

\emph{Points on the unit circle.} Let $\Gamma$ be the countable subgroup of the circle
generated by the remaining points $w$ with $m_w\ne0$, with the discrete topology. Its dual $G$
is compact and metrizable. Write $g_n\in G$ for the character $w\mapsto w^n$. On each window
take nonnegative triangular weights $t_n$ of total mass one, obtained by normalizing
$\min(i+1,H-2-i)$, $0\le i\le H-3$, and zero at the last two positions. They satisfy
$\sup t_n=O(1/H)$ and $\sup\abs{t_{n+1}-t_n}=O(1/H^2)$. The measures
$\sum_nt_nd_n\delta_{g_n}$ on $G$ have total variation at most $O(C)$, so a subsequence
converges weak-star to a finite measure $\nu$. For $z\in\Gamma$, Abel summation shows that
the $t_n$-weighted average of $(w/z)^n$ tends to zero for $w\ne z$, the total variation of the
weights being $O(1/H)$. Since the $c_w$ are absolutely summable, the weighted average of
$a_nz^{-n}$ tends to $c_z$, with $c_z=0$ when $m_z=0$. Shifting the indices in the second
difference gives
\[
 \sum_nt_na_nz^{-n}=-(z-1)^2\sum_nt_nd_nz^{-n}+O\Bigl(H^{-2}\sum_n\abs{d_n}\Bigr),
\]
and the error is $O(C/H)$. Hence $\hat\nu(z)=-m_z$ for $z\ne1$. Subtracting $\hat\nu(1)$
times the Haar probability makes the coefficient at $1$ vanish.

The resulting Fourier transform is bounded and integer-valued. Choose a polynomial that
vanishes at $0$ and takes the value $1$ at every nonzero integer of its finite range.
Applying the corresponding convolution polynomial to the measure gives a finite measure whose
Fourier transform is the indicator of $D=\{w:m_w\ne0\}$. By the idempotent theorem of Cohen,
Helson and Rudin \cite{GreenSanders2008Quantitative}, $D$ lies in the coset ring of $\Gamma$.

An infinite member of the coset ring of a subgroup of the circle is dense in the circle. To
see this, write a Boolean combination of finitely many cosets as a union of atoms. The
positive cosets of an atom intersect in a single coset of a subgroup $T$. Split it into
finitely many smaller cosets so as to settle all excluded cosets whose subgroups meet $T$ in
finite index. Each surviving piece is a coset of a subgroup $T_0$ with finitely many cosets of
subgroups of infinite index removed. If $T_0$ is infinite, it is dense in the circle. For a
translation-invariant finitely additive probability on the abelian group $T_0$, every coset
of infinite index has measure zero, by disjoint translates, and the trace on $T_0$ of every
nonempty open arc has positive measure, since finitely many translates cover $T_0$. So each
surviving infinite piece is dense. But the only possible accumulation point of $D$ is $1$,
because $w_\rho\to1$ as the height grows. Thus $D$ is finite.

\emph{Finitely many differing zeros.} The completed functions $\xi_K,\xi_L$ are entire of
order one, and Hadamard factorization now gives $\xi_K/\xi_L=e^{\alpha+\beta s}P(s)/Q(s)$ with
polynomials $P,Q$. Put $Z=\zeta_K/\zeta_L$. The gamma factors of a number field are products
of $\Gamma_\R(s)$ and $\Gamma_\R(s+1)$, and $\Gamma_\R(s+2)/\Gamma_\R(s)=s/(2\pi)$, while the
discriminant factor contributes a constant under the shift $s\mapsto s+2$. Hence
$Z(s+2)/Z(s)$ is a rational function. On the positive real axis $\zeta_K(s)=1+O_K(2^{-s})$,
and likewise for $L$, so $Z(s+2)/Z(s)=1+O(2^{-s})$. A rational function with this decay
towards $1$ is identically $1$. For real $s>1$ this gives $Z(s)=Z(s+2m)\to1$, so
$\zeta_K=\zeta_L$.
\end{proof}

The theorem identifies a zeta function up to equality, so it does not separate nonisomorphic
arithmetically equivalent fields. It is a statement about infinitely many samples and gives
no effective window length for a given pair of fields.

\subsection{Open questions}\label{ana:sec:open}

\begin{question}[Arithmetic inputs]\label{ana:q:input}
No hypothesis of the criteria of this section is established for a function outside the
scope of \cref{src:thm:dgrh}. Is there such a function, for instance $\zeta_K$ with
$K=\Q(\sqrt[3]{2})$, for which one of the weaker inputs can be proved by a comparison of
covolumes: a bound for the positive part in mean (\cref{ana:cor:forms}), with losses
controlled only on average as in \cref{ana:prop:rankloss}; a signed moving-window bound
(\cref{ana:thm:moving}); a pair of filtered bounds (\cref{ana:thm:joint}); a finite
difference (\cref{ana:thm:fd}); or a bound on a sparse grid
(\cref{ana:thm:powers,ana:cor:finite})? The specific targets are stated in
\cref{sec:grp,sec:hec}.
\end{question}

\begin{question}[Uniform windows]\label{ana:q:window}
The constant $L$ of \cref{ana:thm:window} depends on the family through the radius of a disc
of holomorphy. Can it be chosen uniformly for a family of errors attached to fields of
increasing conductor, for example for increasing blocks of Kummer fields? Uniformity is not
needed for a criterion about finitely many fixed fields, but it is needed for quantitative
error terms and for any criterion about an infinite block.
\end{question}

\begin{question}[Thin twist families beyond Dirichlet characters]\label{ana:q:twist}
\Cref{ana:thm:sparse} rests on nonnegative weights, on the residue formula
\eqref{ana:eq:srcres} and on the zero-density estimate \eqref{ana:eq:zerodensity} for the
twisting family. Is there a family of twists of a nonabelian $L$-function, for instance of the
two-dimensional factor of the splitting field of $X^3-2$, with a zero-density estimate that
permits the same selection?
\end{question}

\begin{question}[A prescribed base on thin grids]\label{ana:q:base}
\Cref{ana:prop:generic} holds for a generic base, and \cref{ana:thm:blocks} for every base on
the particular set $\anaE$. For a prescribed base such as $b=2$ and the grid
$x=b^{\lambda_n}$ with $\lambda_n=\exp(n/\log(n+2))$, do conclusions (i) and (ii) of
\cref{ana:thm:powers} hold?
\end{question}

\begin{question}[Unattained suprema]\label{ana:q:sup}
Let $K$ be a number field for which the real parts of the zeros off the critical line approach
$\Theta_K>1/2$ without attaining it. Is $\sqrt x\,R_K(x)$ unbounded above on the grid
$x=b^e$, $e\in\anaE$? \Cref{ana:thm:blocks} decides this only when $\Theta_K$ is attained,
whereas the Li coefficients on the same index set detect every failure
(\cref{ana:thm:li}).
\end{question}

\begin{question}[Simple zeros]\label{ana:q:simple}
Prove that the sampled mean square in \cref{ana:thm:simple}(iii) is at most the displayed
constant, that is, $D=0$. A multiple zero at very large height would contribute less than the
bound $1.67\cdot10^{-34}$, so no finite computation of zeros decides this.
\end{question}

\begin{question}[The lower limit of the source remainder]\label{ana:q:parity}
Does the parity condition of \cref{ana:prop:parity} hold for the ordinates of $\zeta$ and of
$L(\cdot,\chi_d)$? If it does, $\liminf_{x\to\infty}\sqrt x\,R_d(x)=-B_d-B_0$ exactly. In
general, where in the interval of \cref{ana:thm:excursion} does the lower limit lie?
\end{question}

\begin{question}[Effective rigidity]\label{ana:q:rigidity}
Give an effective window length in \cref{ana:thm:rigidity}: for two given number fields with
different zeta functions, bound the length $H$ beyond which every window has mean
$\abs{d_n}$ above a given size. This requires quantitative control of ordinates, of their
separation, and of the tails as the fields vary.
\end{question}
